\documentclass[aspectratio=169]{beamer}\usepackage[]{graphicx}\usepackage[]{xcolor}
% maxwidth is the original width if it is less than linewidth
% otherwise use linewidth (to make sure the graphics do not exceed the margin)
\makeatletter
\def\maxwidth{ %
  \ifdim\Gin@nat@width>\linewidth
    \linewidth
  \else
    \Gin@nat@width
  \fi
}
\makeatother

\definecolor{fgcolor}{rgb}{0.345, 0.345, 0.345}
\newcommand{\hlnum}[1]{\textcolor[rgb]{0.686,0.059,0.569}{#1}}%
\newcommand{\hlsng}[1]{\textcolor[rgb]{0.192,0.494,0.8}{#1}}%
\newcommand{\hlcom}[1]{\textcolor[rgb]{0.678,0.584,0.686}{\textit{#1}}}%
\newcommand{\hlopt}[1]{\textcolor[rgb]{0,0,0}{#1}}%
\newcommand{\hldef}[1]{\textcolor[rgb]{0.345,0.345,0.345}{#1}}%
\newcommand{\hlkwa}[1]{\textcolor[rgb]{0.161,0.373,0.58}{\textbf{#1}}}%
\newcommand{\hlkwb}[1]{\textcolor[rgb]{0.69,0.353,0.396}{#1}}%
\newcommand{\hlkwc}[1]{\textcolor[rgb]{0.333,0.667,0.333}{#1}}%
\newcommand{\hlkwd}[1]{\textcolor[rgb]{0.737,0.353,0.396}{\textbf{#1}}}%
\let\hlipl\hlkwb

\usepackage{framed}
\makeatletter
\newenvironment{kframe}{%
 \def\at@end@of@kframe{}%
 \ifinner\ifhmode%
  \def\at@end@of@kframe{\end{minipage}}%
  \begin{minipage}{\columnwidth}%
 \fi\fi%
 \def\FrameCommand##1{\hskip\@totalleftmargin \hskip-\fboxsep
 \colorbox{shadecolor}{##1}\hskip-\fboxsep
     % There is no \\@totalrightmargin, so:
     \hskip-\linewidth \hskip-\@totalleftmargin \hskip\columnwidth}%
 \MakeFramed {\advance\hsize-\width
   \@totalleftmargin\z@ \linewidth\hsize
   \@setminipage}}%
 {\par\unskip\endMakeFramed%
 \at@end@of@kframe}
\makeatother

\definecolor{shadecolor}{rgb}{.97, .97, .97}
\definecolor{messagecolor}{rgb}{0, 0, 0}
\definecolor{warningcolor}{rgb}{1, 0, 1}
\definecolor{errorcolor}{rgb}{1, 0, 0}
\newenvironment{knitrout}{}{} % an empty environment to be redefined in TeX

\usepackage{alltt}

% Set lecture number for later use


% Part common to all the lectures
\subtitle{MATH 2740 -- Mathematics of Data Science -- Unit U10}
\author{\texorpdfstring{Julien Arino\newline\url{julien.arino@umanitoba.ca}}{Julien Arino}}
\institute{Department of Mathematics @ University of Manitoba}

% Title of the lecture
\title{Symmetric graph matrices and spectral clustering}



\usetheme{default}
% Slide setup, colour independent

\usepackage{amsmath,amssymb,amsthm}
\usepackage[utf8]{inputenc}
\usepackage{colortbl}
\usepackage{bm}
\usepackage{xcolor}
\usepackage{dsfont}
\usepackage{setspace}
% To use \ding{234} and the like
\usepackage{pifont}
% To cross reference between slide files
\usepackage{zref-xr,zref-user}
% Use something like
% \zexternaldocument{fileI}
% in the tex files. And cite using \zref instead of \ref

% Cross-reference system - see CROSS-REFERENCE-SETUP.md for manual setup instructions
\usepackage{booktabs}
\usepackage{marvosym}
\usepackage{cancel}
%\usepackage{transparent}
% Make doi clickable in the bibliography?
\usepackage{doi}

\usepackage[T1]{fontenc}

\usepackage{longtable}

% For heavier titles
\usepackage{helvet} % Enables Helvetica font family


% Fields and the like
\def\IC{\mathbb{C}}
\def\IE{\mathbb{E}}
\def\IF{\mathbb{F}}
\def\II{\mathbb{I}}
\def\IJ{\mathbb{J}}
\def\IK{\mathbb{K}}
\def\IM{\mathbb{M}}
\def\IN{\mathbb{N}}
\def\IP{\mathbb{P}}
\def\IR{\mathbb{R}}
\newcommand{\IRplus}{\mathbb{R}_{\ge 0}}
\def\IZ{\mathbb{Z}}
\def\11{\mathds{1}}


% Bold lowercase
\def\ba{\bm{a}}
\def\bb{\bm{b}}
\def\bc{\bm{c}}
\def\bd{\bm{d}}
\def\be{\bm{e}}
\def\bf{\bm{f}}
\def\bg{\bm{g}}
\def\bh{\bm{h}}
\def\bi{\bm{i}}
\def\bj{\bm{j}}
\def\bk{\bm{k}}
\def\bn{\bm{n}}
\def\bp{\bm{p}}
\def\br{\bm{r}}
\def\bs{\bm{s}}
\def\bu{\bm{u}}
\def\bv{\bm{v}}
\def\bw{\bm{w}}
\def\bx{\bm{x}}
\def\by{\bm{y}}
\def\bz{\bm{z}}
\newcommand{\vect}[1]{\bm{#1}}

% Bold capitals
\def\bB{\bm{B}}
\def\bD{\bm{D}}
\def\bE{\bm{E}}
\def\bF{\bm{F}}
\def\bG{\bm{G}}
\def\bI{\bm{I}}
\def\bL{\bm{L}}
\def\bN{\bm{N}}
\def\bP{\bm{P}}
\def\bR{\bm{R}}
\def\bS{\bm{S}}
\def\bT{\bm{T}}
\def\bX{\bm{X}}

% Bold numbers
\def\b0{\bm{0}}

% Bold greek
\bmdefine{\bmu}{\bm{\mu}}
\def\bphi{\bm{\phi}}
\def\bvarphi{\bm{\varphi}}
\def\bPi{\bm{\Pi}}
\def\bGamma{\bm{\Gamma}}

% Bold red sentence
\def\boldred#1{{\color{red}\textbf{#1}}}
\def\defword#1{{\color{orange}\textbf{#1}}}

% Caligraphic letters
\def\A{\mathcal{A}}
\def\B{\mathcal{B}}
\def\C{\mathcal{C}}
\def\D{\mathcal{D}}
\def\E{\mathcal{E}}
\def\F{\mathcal{F}}
\def\G{\mathcal{G}}
\def\H{\mathcal{H}}
\def\I{\mathcal{I}}
\def\L{\mathcal{L}}
\def\M{\mathcal{M}}
\def\N{\mathcal{N}}
\def\P{\mathcal{P}}
\def\R{\mathcal{R}}
\def\S{\mathcal{S}}
\def\T{\mathcal{T}}
\def\U{\mathcal{U}}
\def\V{\mathcal{V}}

% Adding space for prime (') where needed
\def\pprime{\,'}
% Adding space for star (\star) where needed
\def\pstar{{\,\star}}

% tt font for code
\def\code#1{{\tt #1}}

% i.e., e.g.
\def\eg{\emph{e.g.}}
\def\ie{\emph{i.e.}}


% Operators and special symbols
\def\nbOne{{\mathchoice {\rm 1\mskip-4mu l} {\rm 1\mskip-4mu l}
{\rm 1\mskip-4.5mu l} {\rm 1\mskip-5mu l}}}
\def\cov{\ensuremath{\mathsf{cov}}}
\def\Var{\ensuremath{\mathsf{Var}\ }}
\def\Im{\textrm{Im}\;}
\def\Re{\textrm{Re}\;}
\def\det{\ensuremath{\mathsf{det}}}
\def\diag{\ensuremath{\mathsf{diag}}}
\def\nullspace{\ensuremath{\mathsf{null}}}
\def\nullity{\ensuremath{\mathsf{nullity}}}
\def\rank{\ensuremath{\mathsf{rank}}}
\def\range{\ensuremath{\mathsf{range}}}
\def\sgn{\ensuremath{\mathsf{sgn}}}
\def\Span{\ensuremath{\mathsf{span}}}
\def\tr{\ensuremath{\mathsf{tr}}}
\def\imply{$\Rightarrow$}
\def\restrictTo#1#2{\left.#1\right|_{#2}}
\newcommand{\parallelsum}{\mathbin{\!/\mkern-5mu/\!}}
\def\dsum{\mathop{\displaystyle \sum }}%
\def\dind#1#2{_{\substack{#1\\ #2}}}

\newcommand{\Qmatrix}[1]{%
  \begin{pmatrix}#1\end{pmatrix}%
}

\DeclareMathOperator{\GL}{GL}
\DeclareMathOperator{\Rel}{Re}
\def\Nt#1{\left|\!\left|\!\left|#1\right|\!\right|\!\right|}
\newcommand{\tripbar}{|\! |\! |}



% The beamer bullet (in base colour)
\def\bbullet{\leavevmode\usebeamertemplate{itemize item}\ }

% Theorems and the like
\newtheorem{importantdefinition}[theorem]{Definition}
\newtheorem{proposition}[theorem]{Proposition}
\newtheorem{property}[theorem]{Property}
\newtheorem{importantproperty}[theorem]{Property}
\newtheorem{importanttheorem}[theorem]{Theorem}
%\newtheorem{lemma}[theorem]{Lemma}
%\newtheorem{corollary}[theorem]{Corollary}
\newtheorem{remark}[theorem]{Remark}
\setbeamertemplate{theorems}[numbered]
%\setbeamertemplate{theorems}[ams style]

%
%\usecolortheme{orchid}
%\usecolortheme{orchid}

\def\red{\color[rgb]{1,0,0}}
\def\blue{\color[rgb]{0,0,1}}
\def\green{\color[rgb]{0,1,0}}

% Fix skipping lines after items in the bibliography
\setbeamertemplate{bibliography entry title}{}
\setbeamertemplate{bibliography entry location}{}
\setbeamertemplate{bibliography entry note}{}

% Get rid of navigation stuff
\setbeamertemplate{navigation symbols}{}

% Set footline/header line
\setbeamertemplate{footline}
{%
\quad p. \insertpagenumber \quad--\quad \insertsection\vskip2pt
}
% \setbeamertemplate{headline}
% {%
% \quad\insertsection\hfill p. \insertpagenumber\quad\mbox{}\vskip2pt
% }


\makeatletter
\newlength\beamerleftmargin
\setlength\beamerleftmargin{\Gm@lmargin}
\makeatother

% Colours for special pages
\def\extraContent{yellow!20}


%%%%%%%%%%%%%%%%%
\usepackage{tikz}
\usetikzlibrary{shapes,arrows}
\usetikzlibrary{positioning}
\usetikzlibrary{shapes.symbols,shapes.callouts,patterns}
\usetikzlibrary{calc,fit}
\usetikzlibrary{backgrounds}
\usetikzlibrary{decorations.pathmorphing,fit,petri}
\usetikzlibrary{automata}
\usetikzlibrary{fadings}
\usetikzlibrary{patterns,hobby}
\usetikzlibrary{backgrounds,fit,petri}
\usetikzlibrary{tikzmark}

\usepackage{pgfplots}
\pgfplotsset{compat=1.6}
\pgfplotsset{ticks=none}

\usetikzlibrary{decorations.markings}
\usetikzlibrary{arrows.meta}
\tikzset{>=stealth}

% For tikz
\tikzstyle{cloud} = [draw, ellipse,fill=red!20, node distance=0.87cm,
minimum height=2em]
\tikzstyle{line} = [draw, -latex']


%%% For max frame images
\newenvironment{changemargin}[2]{%
\begin{list}{}{%
\setlength{\topsep}{0pt}%
\setlength{\leftmargin}{#1}%
\setlength{\rightmargin}{#2}%
\setlength{\listparindent}{\parindent}%
\setlength{\itemindent}{\parindent}%
\setlength{\parsep}{\parskip}%
}%
\item[]}{\end{list}}


% Make one image take up the entire slide content area in beamer,.:
% centered/centred full-screen image, with title:
% This uses the whole screen except for the 1cm border around it
% all. 128x96mm
\newcommand{\titledFrameImage}[2]{
\begin{frame}{#1}
%\begin{changemargin}{-1cm}{-1cm}
\begin{center}
\includegraphics[width=108mm,height=\textheight,keepaspectratio]{#2}
\end{center}
%\end{changemargin}
\end{frame}
}

% Make one image take up the entire slide content area in beamer.:
% centered/centred full-screen image, no title:
% This uses the whole screen except for the 1cm border around it
% all. 128x96mm
\newcommand{\plainFrameImage}[1]{
\begin{frame}[plain]
%\begin{changemargin}{-1cm}{-1cm}
\begin{center}
\includegraphics[width=108mm,height=76mm,keepaspectratio]{#1}
\end{center}
%\end{changemargin}
\end{frame}
}

% Make one image take up the entire slide area, including borders, in beamer.:
% centered/centred full-screen image, no title:
% This uses the entire whole screen
\newcommand{\maxFrameImage}[1]{
\begin{frame}[plain]
\begin{changemargin}{-1cm}{-1cm}
\begin{center}
\includegraphics[width=\paperwidth,height=\paperheight,keepaspectratio]
{#1}
\end{center}
\end{changemargin}
\end{frame}
}

% This uses the entire whole screen (to include in frame)
\newcommand{\maxFrameImageNoFrame}[1]{
\begin{changemargin}{-1cm}{-1cm}
\begin{center}
\includegraphics[width=\paperwidth,height=0.99\paperheight,keepaspectratio]
{#1}
\end{center}
\end{changemargin}
}

% Make one image take up the entire slide area, including borders, in beamer.:
% centered/centred full-screen image, no title:
% This uses the entire whole screen
\newcommand{\maxFrameImageColor}[2]{
\begin{frame}[plain]
\setbeamercolor{normal text}{bg=#2!20}
\begin{changemargin}{-1cm}{-1cm}
\begin{center}
\includegraphics[width=\paperwidth,height=\paperheight,keepaspectratio]
{#1}
\end{center}
\end{changemargin}
\end{frame}
}


\usepackage{tikz}
\usetikzlibrary{patterns,hobby,matrix}
\usepackage{pgfplots}
\pgfplotsset{compat=1.6}
\pgfplotsset{ticks=none}

\usetikzlibrary{backgrounds}
\usetikzlibrary{decorations.markings}
\usetikzlibrary{arrows.meta}
\tikzset{>=stealth}

\tikzset{
  clockwise arrows/.style={
    postaction={
      decorate,
      decoration={
        markings,
        mark=between positions 0.1 and 0.9 step 40pt with {\arrow{>}},
   }}}}


% New integrated section command: creates section and section slide
\newcommand{\Ssection}[2]{
\section{#1}
\begin{frame}[noframenumbering,plain]
  \begin{tikzpicture}[remember picture,overlay]
    \node[above right,inner sep=0pt,opacity=0.2] at (current page.south west)
    {
        \includegraphics[height=\paperheight,width=\paperwidth]{#2}
    };
  \end{tikzpicture}
  \setbeamercolor{section in toc}{fg=section_page_list_colour}
  \setbeamerfont{section in toc}{size=\Large,series=\bfseries}
  \setbeamertemplate{section in toc shaded}[default][60]
  \tableofcontents[
    currentsection,
    sectionstyle=show/shaded,
    subsectionstyle=show/hide/hide,
    subsubsectionstyle=hide/hide/hide]
\end{frame}
\addtocounter{page}{-1}
}

% New integrated section command with subsections: creates section and section slide showing subsections
\newcommand{\SsectionWithSubs}[2]{
\section{#1}
\begin{frame}[noframenumbering,plain]
  \begin{tikzpicture}[remember picture,overlay]
    \node[above right,inner sep=0pt,opacity=0.2] at (current page.south west)
    {
        \includegraphics[height=\paperheight,width=\paperwidth]{#2}
    };
  \end{tikzpicture}
  \setbeamercolor{section in toc}{fg=section_page_list_colour}
  \setbeamerfont{section in toc}{size=\Large,series=\bfseries}
  \setbeamertemplate{section in toc shaded}[default][60]
  \tableofcontents[
    currentsection,
    sectionstyle=show/hide,
    subsectionstyle=show/show/hide,
    subsubsectionstyle=hide/hide/hide]
\end{frame}
\addtocounter{page}{-1}
}

% New integrated subsection command: creates subsection and subsection slide
\newcommand{\Ssubsection}[2]{
\subsection{#1}
\begin{frame}[noframenumbering,plain]
  \begin{tikzpicture}[remember picture,overlay]
    \node[above right,inner sep=0pt,opacity=0.2] at (current page.south west)
    {
        \includegraphics[height=\paperheight,width=\paperwidth]{#2}
    };
  \end{tikzpicture}
  \setbeamercolor{section in toc}{fg=subsection_page_list_colour}
  \setbeamerfont{section in toc}{size=\Large,series=\bfseries}
  \setbeamertemplate{section in toc shaded}[default][60]
  \setbeamerfont{subsection in toc}{series=\bfseries}
  \setbeamertemplate{subsection in toc shaded}[default][50]
  \tableofcontents[
    currentsection,
    sectionstyle=show/hide,
    subsectionstyle=show/shaded/hide,
    subsubsectionstyle=hide/hide/hide]
\end{frame}
\addtocounter{page}{-1}
}

% New integrated subsubsection command: creates subsubsection and subsubsection slide
\newcommand{\Ssubsubsection}[2]{
\subsubsection{#1}
\begin{frame}[noframenumbering,plain]
  \begin{tikzpicture}[remember picture,overlay]
    \node[above right,inner sep=0pt,opacity=0.2] at (current page.south west)
    {
        \includegraphics[height=\paperheight,width=\paperwidth]{#2}
    };
  \end{tikzpicture}
  \setbeamercolor{section in toc}{fg=subsub_header_section}
  \setbeamerfont{section in toc}{size=\Large,series=\bfseries}
  \setbeamertemplate{section in toc shaded}[default][60]
  \setbeamerfont{subsection in toc}{series=\bfseries}
  \setbeamertemplate{subsection in toc shaded}[default][50]
  \setbeamertemplate{subsubsection in toc shaded}[default][50]
  \tableofcontents[
    currentsection,
    sectionstyle=show/hide,
    subsectionstyle=show/hide/hide,
    subsubsectionstyle=show/shaded/hide]
\end{frame}
\addtocounter{page}{-1}
}

% Legacy commands (kept for backward compatibility)
% Beginning of a section
\newcommand{\newSectionSlide}[1]{
\begin{frame}[noframenumbering,plain]
  \begin{tikzpicture}[remember picture,overlay]
    \node[above right,inner sep=0pt,opacity=0.2] at (current page.south west)
    {
        \includegraphics[height=\paperheight,width=\paperwidth]{#1}
    };
  \end{tikzpicture}
  \setbeamercolor{section in toc}{fg=section_page_list_colour}
  \setbeamerfont{section in toc}{size=\Large,series=\bfseries}
  \setbeamertemplate{section in toc shaded}[default][60]
  \tableofcontents[
    currentsection,
    sectionstyle=show/shaded,
    subsectionstyle=show/hide/hide,
    subsubsectionstyle=hide/hide/hide]
\end{frame}
\addtocounter{page}{-1}
}

% Beginning of a section in which we also show subsections
\newcommand{\newSectionWithSubsSlide}[1]{
	\begin{frame}[noframenumbering,plain]
		\begin{tikzpicture}[remember picture,overlay]
			\node[above right,inner sep=0pt,opacity=0.2] at (current page.south west)
			{
				\includegraphics[height=\paperheight,width=\paperwidth]{#1}
			};
		\end{tikzpicture}
		\setbeamercolor{section in toc}{fg=section_page_list_colour}
		\setbeamerfont{section in toc}{size=\Large,series=\bfseries}
		\setbeamertemplate{section in toc shaded}[default][60]
		\tableofcontents[
		currentsection,
		sectionstyle=show/hide,
		subsectionstyle=show/show/hide,
		subsubsectionstyle=hide/hide/hide]
	\end{frame}
	\addtocounter{page}{-1}
}

% Beginning of a subsection
\newcommand{\newSubSectionSlide}[1]{
\begin{frame}[noframenumbering,plain]
  \begin{tikzpicture}[remember picture,overlay]
    \node[above right,inner sep=0pt,opacity=0.2] at (current page.south west)
    {
        \includegraphics[height=\paperheight,width=\paperwidth]{#1}
    };
  \end{tikzpicture}
  \setbeamercolor{section in toc}{fg=subsection_page_list_colour}
  \setbeamerfont{section in toc}{size=\Large,series=\bfseries}
  \setbeamertemplate{section in toc shaded}[default][60]
  \setbeamerfont{subsection in toc}{series=\bfseries}
  \setbeamertemplate{subsection in toc shaded}[default][50]
  \tableofcontents[
    currentsection,
    sectionstyle=show/hide,
    subsectionstyle=show/shaded/hide,
    subsubsectionstyle=hide/hide/hide]
\end{frame}
\addtocounter{page}{-1}
}

% Beginning of a subsubsection
\newcommand{\newSubSubSectionSlide}[1]{
\begin{frame}[noframenumbering,plain]
  \begin{tikzpicture}[remember picture,overlay]
    \node[above right,inner sep=0pt,opacity=0.2] at (current page.south west)
    {
        \includegraphics[height=\paperheight,width=\paperwidth]{#1}
    };
  \end{tikzpicture}
  \setbeamercolor{section in toc}{fg=subsub_header_section}
  \setbeamerfont{section in toc}{size=\Large,series=\bfseries}
  \setbeamertemplate{section in toc shaded}[default][60]
  \setbeamerfont{subsection in toc}{series=\bfseries}
  \setbeamertemplate{subsection in toc shaded}[default][50]
  \setbeamertemplate{subsubsection in toc shaded}[default][50]
  \tableofcontents[
    currentsection,
    sectionstyle=show/hide,
    subsectionstyle=show/hide/hide,
    subsubsectionstyle=show/shaded/hide]
\end{frame}
\addtocounter{page}{-1}
}


   %%%%%%%%%%%
% To have links to parts in the outline
\makeatletter
\AtBeginPart{%
  \addtocontents{toc}{\protect\beamer@partintoc{\the\c@part}{\beamer@partnameshort}{\the\c@page}}%
}
%% number, shortname, page.
\providecommand\beamer@partintoc[3]{%
  \ifnum\c@tocdepth=-1\relax
    % requesting onlyparts.
    \makebox[6em]{Part #1:} \textcolor{green!30!blue}{\hyperlink{#2}{#2}}
    \par
  \fi
}
\define@key{beamertoc}{onlyparts}[]{%
  \c@tocdepth=-1\relax
}
\makeatother%

\newcommand{\nameofthepart}{}
\newcommand{\nupart}[1]%
    {   \part{#1}%
        \renewcommand{\nameofthepart}{#1}%
        {
          \setbeamercolor{background canvas}{bg=orange!50}
          \begin{frame}{#1}%\partpage 
          \hypertarget{\nameofthepart}{}\tableofcontents%
          \end{frame}
        }
    }

% This command creates a title page using TikZ only
\newcommand{\tikztitlepage}[1]{%
\begin{frame}[plain,noframenumbering]
  \begin{tikzpicture}[remember picture,overlay]
    % Background image
    \node[above right,inner sep=0pt,opacity=0.1] 
      at (current page.south west) 
      {\includegraphics[width=\paperwidth,height=\paperheight]{#1}};

    % University logo
    \node[anchor=north east, inner sep=5pt, opacity=0.9] 
      at (current page.north east)
      {\includegraphics[width=0.2\textwidth]{FIGS-slides-admin/UM-logo-horizontal-CMYK.png}};
    
    % Title
    \node[anchor=center, align=center, 
          font=\fontsize{13}{15}\bfseries\color{UMbrown}, 
          text width=0.9\textwidth] 
          at ([yshift=2cm]current page.center)
          {\inserttitle};

      % Authors
      \node[anchor=center, align=center,
        font=\fontsize{10}{12}\bfseries\color{UMbrown},
        text width=0.7\textwidth]
        at ([yshift=0.8cm]current page.center)
        {\insertauthor};

      % Affiliation
      \node[anchor=north, align=center,
        font=\fontsize{9}{11}\color{UMbrown},
        text width=0.7\textwidth]
        at ([yshift=-0.2cm]current page.center)
        {\insertaffiliation};      
    % Date
    \node[anchor=north, align=center, 
          font=\fontsize{12}{16}\bfseries\color{UMbrown},
          text width=0.7\textwidth] 
          at ([yshift=0.2cm]current page.center)
          {\insertdate};

    % Land acknowledgement
    \node[anchor=south, align=justify, 
          font=\footnotesize, text=black, 
          text width=1.1\textwidth] 
          at ([yshift=0.5cm]current page.south)
          {The University of Manitoba campuses are located on original lands of Anishinaabeg, Ininew, Anisininew, Dakota and Dene peoples, and on the National Homeland of the Red River Métis.\\
          We respect the Treaties that were made on these territories, we acknowledge the harms and mistakes of the past, and we dedicate ourselves to move forward in partnership with Indigenous communities in a spirit of Reconciliation and collaboration.};
  \end{tikzpicture}
  \addtocounter{page}{-1}
\end{frame}
}
% The title page with figure
% \newcommand{\titlepagewithfigure}[1]{%
%   \begin{frame}[noframenumbering,plain]
%     \begin{tikzpicture}[remember picture,overlay]
%       \node[above right,inner sep=0pt,opacity=0.1] at (current page.south west)
%       {
%           \includegraphics[height=\paperheight,width=\paperwidth]{#1}
%       };
%       \node[anchor=north east,
%       inner sep=5pt,
%       opacity=0.9] at (current page.north east)
%       {
%           \includegraphics[width=0.2\textwidth]{FIGS-slides-admin/UM-logo-horizontal-CMYK.png}
%       };
%       \node[anchor=south, 
%       align=justify, 
%       text=black, 
%       text width=1.1\textwidth,
%       font=\footnotesize]  (land_acknowledgement)
%       at (current page.south) 
%       {The University of Manitoba campuses are located on original lands of Anishinaabeg, Ininew, Anisininew, Dakota and Dene peoples, and on the National Homeland of the Red River Métis.
%       We respect the Treaties that were made on these territories, we acknowledge the harms and mistakes of the past, and we dedicate ourselves to move forward in partnership with Indigenous communities in a spirit of Reconciliation and collaboration.};  
%       % \node[align=center, anchor=south,
%       % above=0.5cm of land_acknowledgement,
%       % text=black,
%       % font=\bfseries] {\insertdate};
%   \end{tikzpicture}
%   \setbeamercolor{title}{fg=title_page_title_colour}
%   \setbeamerfont{title}{size=\Large,series=\bfseries}
%   \setbeamercolor{author}{fg=title_page_author_colour}
%   \setbeamerfont{author}{size=\large,series=\bfseries}
%   \setbeamercolor{institute}{fg=title_page_institute_colour}
%   \setbeamerfont{institute}{size=\large,series=\bfseries}
%   \setbeamercolor{date}{fg=title_page_date_colour}
%   \setbeamerfont{date}{series=\bfseries}
% 	\titlepage
% \end{frame}
% \addtocounter{page}{-1}
% }

\newcommand{\titlepagewithfigure}[1]{%
  \begin{frame}[noframenumbering,plain]
    \begin{tikzpicture}[remember picture,overlay]
      \node[above right,inner sep=0pt,opacity=0.1] at (current page.south west)
      {
          \includegraphics[height=\paperheight,width=\paperwidth]{#1}
      };
      \node[anchor=north east,
      inner sep=5pt,
      opacity=0.9] at (current page.north east)
      {
          \includegraphics[width=0.2\textwidth]{FIGS-slides-admin/UM-logo-horizontal-CMYK.png}
      };
      \node[anchor=south, 
      align=justify, 
      text=black, 
      text width=1.1\textwidth,
      font=\footnotesize]  (land_acknowledgement)
      at (current page.south) 
      {The University of Manitoba campuses are located on original lands of Anishinaabeg, Ininew, Anisininew, Dakota and Dene peoples, and on the National Homeland of the Red River Métis.
      We respect the Treaties that were made on these territories, we acknowledge the harms and mistakes of the past, and we dedicate ourselves to move forward in partnership with Indigenous communities in a spirit of Reconciliation and collaboration.};  
      % \node[align=center, anchor=south,
      % above=0.5cm of land_acknowledgement,
      % text=black,
      % font=\bfseries] {\insertdate};
  \end{tikzpicture}
  \setbeamercolor{title}{fg=title_page_title_colour}
  \setbeamerfont{title}{size=\Large,series=\bfseries,family=\usefont{T1}{phv}{b}{n}}
  \setbeamercolor{author}{fg=title_page_author_colour}
  \setbeamerfont{author}{size=\large,series=\bfseries,family=\usefont{T1}{phv}{b}{n}}
  \setbeamercolor{institute}{fg=title_page_institute_colour}
  \setbeamerfont{institute}{size=\large,series=\bfseries,family=\usefont{T1}{phv}{b}{n}}
  \setbeamercolor{date}{fg=title_page_date_colour}
  \setbeamerfont{date}{series=\bfseries,family=\usefont{T1}{phv}{b}{n}}
	\titlepage
\end{frame}
\addtocounter{page}{-1}
}
% The outline page, with figure
% \newcommand{\outlinepage}[1]{%
% \begin{frame}[noframenumbering,plain]
%   \begin{tikzpicture}[remember picture,overlay]
%     \node[above right,inner sep=0pt,opacity=0.2] at (current page.south west)
%     {
%         \includegraphics[height=\paperheight,width=\paperwidth]{#1}
%     };
%   \end{tikzpicture}
%   \setbeamercolor{section in toc}{fg=outline_page_list_colour}
%   \setbeamerfont{section in toc}{size=\Large,series=\bfseries,family=\sffamily}
%   \frametitle{\textcolor{outline_page_title_colour}{\LARGE\bfseries Outline}}
%   \tableofcontents[hideallsubsections]
% \end{frame}
% \addtocounter{page}{-1}
% }
% The outline page, with figure
\newcommand{\outlinepage}[1]{%
\begin{frame}[noframenumbering,plain]
  \begin{tikzpicture}[remember picture,overlay]
    \node[above right,inner sep=0pt,opacity=0.2] at (current page.south west)
    {
        \includegraphics[height=\paperheight,width=\paperwidth]{#1}
    };
  \end{tikzpicture}
  \setbeamercolor{section in toc}{fg=outline_page_list_colour}
  % Use Helvetica Bold only for the outline slide TOC
  \setbeamerfont{section in toc}{size=\Large,family=\usefont{T1}{phv}{b}{n}}
  % Use Helvetica Bold for the outline title
  \frametitle{\textcolor{outline_page_title_colour}{\usefont{T1}{phv}{b}{n}\LARGE Outline}}
  \tableofcontents[hideallsubsections]
\end{frame}
\addtocounter{page}{-1}
}


%\let\oldsection\section
%\renewcommand{\section}[2]{\oldsection[#1]\newSectionSlide[#2]}


%%%%%%%%%%%%%%%%%%%%%
% CUSTOM SLIDE BACKGROUNDS
%%%%%%%%%%%%%%%%%%%%%
% Define custom background templates for different colors
\defbeamertemplate*{background canvas}{blue}{%
  \color{blue!15}\vrule width\paperwidth height\paperheight%
}
\defbeamertemplate*{background canvas}{green}{%
  \color{green!15}\vrule width\paperwidth height\paperheight%
}
\defbeamertemplate*{background canvas}{red}{%
  \color{red!15}\vrule width\paperwidth height\paperheight%
}
\defbeamertemplate*{background canvas}{yellow}{%
  \color{yellow!20}\vrule width\paperwidth height\paperheight%
}
\defbeamertemplate*{background canvas}{purple}{%
  \color{purple!15}\vrule width\paperwidth height\paperheight%
}
\defbeamertemplate*{background canvas}{orange}{%
  \color{orange!20}\vrule width\paperwidth height\paperheight%
}

% Define keys for the different background options
\makeatletter
\define@key{beamerframe}{blue}[true]{\setbeamertemplate{background canvas}[blue]}
\define@key{beamerframe}{green}[true]{\setbeamertemplate{background canvas}[green]}
\define@key{beamerframe}{red}[true]{\setbeamertemplate{background canvas}[red]}
\define@key{beamerframe}{yellow}[true]{\setbeamertemplate{background canvas}[yellow]}
\define@key{beamerframe}{purple}[true]{\setbeamertemplate{background canvas}[purple]}
\define@key{beamerframe}{orange}[true]{\setbeamertemplate{background canvas}[orange]}
\makeatother

% Reset to normal background for all frames by default
\BeforeBeginEnvironment{frame}{\setbeamertemplate{background canvas}[default]}

\newcommand{\punnett}[2]{
    \begin{center}
    \renewcommand{\arraystretch}{1.5} % Add space to rows
    \begin{tabular}{|c c | c|c|}
        \multicolumn{2}{c}{} & \multicolumn{2}{c}{\textbf{Father}} \\
        \multicolumn{2}{c}{} & #1 \\ \hline
        #2
    \end{tabular}
    \end{center}
}

% Colour definitions for Punnett squares
\usepackage{multirow}
\usepackage{colortbl}
\colorlet{punnettorange}{orange!30}
\colorlet{punnettblack}{black!50}
\colorlet{punnetttortie}{orange!60!black}

\usetikzlibrary{positioning, decorations.pathreplacing, arrows.meta}

% --- TIKZ STYLES FOR CAT DIAGRAMS ---
\tikzset{
    % This is the main placeholder style for the cat images
    catnode/.style={
        draw, 
        rectangle, 
        rounded corners, 
        minimum height=1.5cm, 
        minimum width=2.5cm, 
        align=center, 
        font=\small\bfseries
    },
    % Allele styles (the X and Y at the top/side)
    allele/.style={font=\Large\bfseries},
    female/.style={allele, color=purple!80!black},
    male/.style={allele, color=blue!80!black},
    % --- Placeholder styles for different cats ---
    % YOU CAN REPLACE THE CONTENTS OF THESE NODES WITH YOUR IMAGES
    orange_cat/.style={
        catnode, 
        fill=orange!30, 
        text=black
    },
    black_cat/.style={
        catnode, 
        fill=black!70, 
        text=white
    },
    tortie_cat/.style={
        catnode, 
        fill=orange!50!black, % A mix for tortoiseshell
        text=white
    }
}

\usetikzlibrary{
    arrows.meta, % For nicer arrow heads (e.g., -Latex)
    positioning, % For relative node placement (e.g., above=of)
    automata     % For state diagrams, loops
}
% --- END TIKZ STYLES ---

% Common date for all slides
\date{Fall 2026}
% Counter for random transition slides generated by R
\newcounter{sssectionbg}
\setcounter{sssectionbg}{1}
\newcommand{\SSsection}[1]{
  \Ssection{#1}{FIGS/random_bg_\thesssectionbg.pdf}
  \stepcounter{sssectionbg}
}

% Random transition slide generated by R for the title page
\newcommand{\SStitlepage}{
  \titlepagewithfigure{FIGS/random_title_bg.pdf}
}

% Random transition slide generated by R for the outline page
\newcommand{\SSoutlinepage}{
  \outlinepage{FIGS/random_outline_bg.pdf}
}
\newcommand{\SSsubsection}[1]{
  \Ssubsection{#1}{FIGS/random_bg_\thesssectionbg.pdf}
  \stepcounter{sssectionbg}
}

\newcommand{\SSsubsubsection}[1]{
  \Ssubsubsection{#1}{FIGS/random_bg_\thesssectionbg.pdf}
  \stepcounter{sssectionbg}
}



\usecolortheme{orchid}
%% Listings
\usepackage{listings}
\definecolor{mygreen}{rgb}{0,0.6,0}
\definecolor{mygray}{rgb}{0.5,0.5,0.5}
\definecolor{mymauve}{rgb}{0.58,0,0.82}
\definecolor{mygold}{rgb}{1,0.843,0}
\definecolor{myblue}{rgb}{0.537,0.812,0.941}

\definecolor{mygold2}{RGB}{120,105,22}
\definecolor{mygrey2}{RGB}{50,50,50}

\definecolor{lgreen}{rgb}{0.6,0.9,.6}
\definecolor{lred}{rgb}{1,0.5,.5}

\lstloadlanguages{R}
\lstset{ %
  language=R,
  backgroundcolor=\color{black!05},   % choose the background color
  basicstyle=\footnotesize\ttfamily,        % size of fonts used for the code
  breaklines=true,                 % automatic line breaking only at whitespace
  captionpos=b,                    % sets the caption-position to bottom
  commentstyle=\color{mygreen},    % comment style
  escapeinside={\%*}{*)},          % if you want to add LaTeX within your code
  keywordstyle=\color{red},       % keyword style
  stringstyle=\color{mygold},     % string literal style
  keepspaces=true,
  columns=fullflexible,
  tabsize=4,
}
% Could also do (in lstset)
% basicstyle==\fontfamily{pcr}\footnotesize
\lstdefinelanguage{Renhanced}%
  {keywords={abbreviate,abline,abs,acos,acosh,action,add1,add,%
      aggregate,alias,Alias,alist,all,anova,any,aov,aperm,append,apply,%
      approx,approxfun,apropos,Arg,args,array,arrows,as,asin,asinh,%
      atan,atan2,atanh,attach,attr,attributes,autoload,autoloader,ave,%
      axis,backsolve,barplot,basename,besselI,besselJ,besselK,besselY,%
      beta,binomial,body,box,boxplot,break,browser,bug,builtins,bxp,by,%
      c,C,call,Call,case,cat,category,cbind,ceiling,character,char,%
      charmatch,check,chol,chol2inv,choose,chull,class,close,cm,codes,%
      coef,coefficients,co,col,colnames,colors,colours,commandArgs,%
      comment,complete,complex,conflicts,Conj,contents,contour,%
      contrasts,contr,control,helmert,contrib,convolve,cooks,coords,%
      distance,coplot,cor,cos,cosh,count,fields,cov,covratio,wt,CRAN,%
      create,crossprod,cummax,cummin,cumprod,cumsum,curve,cut,cycle,D,%
      data,dataentry,date,dbeta,dbinom,dcauchy,dchisq,de,debug,%
      debugger,Defunct,default,delay,delete,deltat,demo,de,density,%
      deparse,dependencies,Deprecated,deriv,description,detach,%
      dev2bitmap,dev,cur,deviance,off,prev,,dexp,df,dfbetas,dffits,%
      dgamma,dgeom,dget,dhyper,diag,diff,digamma,dim,dimnames,dir,%
      dirname,dlnorm,dlogis,dnbinom,dnchisq,dnorm,do,dotplot,double,%
      download,dpois,dput,drop,drop1,dsignrank,dt,dummy,dump,dunif,%
      duplicated,dweibull,dwilcox,dyn,edit,eff,effects,eigen,else,%
      emacs,end,environment,env,erase,eval,equal,evalq,example,exists,%
      exit,exp,expand,expression,External,extract,extractAIC,factor,%
      fail,family,fft,file,filled,find,fitted,fivenum,fix,floor,for,%
      For,formals,format,formatC,formula,Fortran,forwardsolve,frame,%
      frequency,ftable,ftable2table,function,gamma,Gamma,gammaCody,%
      gaussian,gc,gcinfo,gctorture,get,getenv,geterrmessage,getOption,%
      getwd,gl,glm,globalenv,gnome,GNOME,graphics,gray,grep,grey,grid,%
      gsub,hasTsp,hat,heat,help,hist,home,hsv,httpclient,I,identify,if,%
      ifelse,Im,image,\%in\%,index,influence,measures,inherits,install,%
      installed,integer,interaction,interactive,Internal,intersect,%
      inverse,invisible,IQR,is,jitter,kappa,kronecker,labels,lapply,%
      layout,lbeta,lchoose,lcm,legend,length,levels,lgamma,library,%
      licence,license,lines,list,lm,load,local,locator,log,log10,log1p,%
      log2,logical,loglin,lower,lowess,ls,lsfit,lsf,ls,machine,Machine,%
      mad,mahalanobis,make,link,margin,match,Math,matlines,mat,matplot,%
      matpoints,matrix,max,mean,median,memory,menu,merge,methods,min,%
      missing,Mod,mode,model,response,mosaicplot,mtext,mvfft,na,nan,%
      names,omit,nargs,nchar,ncol,NCOL,new,next,NextMethod,nextn,%
      nlevels,nlm,noquote,NotYetImplemented,NotYetUsed,nrow,NROW,null,%
      numeric,\%o\%,objects,offset,old,on,Ops,optim,optimise,optimize,%
      options,or,order,ordered,outer,package,packages,page,pairlist,%
      pairs,palette,panel,par,parent,parse,paste,path,pbeta,pbinom,%
      pcauchy,pchisq,pentagamma,persp,pexp,pf,pgamma,pgeom,phyper,pico,%
      pictex,piechart,Platform,plnorm,plogis,plot,pmatch,pmax,pmin,%
      pnbinom,pnchisq,pnorm,points,poisson,poly,polygon,polyroot,pos,%
      postscript,power,ppoints,ppois,predict,preplot,pretty,Primitive,%
      print,prmatrix,proc,prod,profile,proj,prompt,prop,provide,%
      psignrank,ps,pt,ptukey,punif,pweibull,pwilcox,q,qbeta,qbinom,%
      qcauchy,qchisq,qexp,qf,qgamma,qgeom,qhyper,qlnorm,qlogis,qnbinom,%
      qnchisq,qnorm,qpois,qqline,qqnorm,qqplot,qr,Q,qty,qy,qsignrank,%
      qt,qtukey,quantile,quasi,quit,qunif,quote,qweibull,qwilcox,%
      rainbow,range,rank,rbeta,rbind,rbinom,rcauchy,rchisq,Re,read,csv,%
      csv2,fwf,readline,socket,real,Recall,rect,reformulate,regexpr,%
      relevel,remove,rep,repeat,replace,replications,report,require,%
      resid,residuals,restart,return,rev,rexp,rf,rgamma,rgb,rgeom,R,%
      rhyper,rle,rlnorm,rlogis,rm,rnbinom,RNGkind,rnorm,round,row,%
      rownames,rowsum,rpois,rsignrank,rstandard,rstudent,rt,rug,runif,%
      rweibull,rwilcox,sample,sapply,save,scale,scan,scan,screen,sd,se,%
      search,searchpaths,segments,seq,sequence,setdiff,setequal,set,%
      setwd,show,sign,signif,sin,single,sinh,sink,solve,sort,source,%
      spline,splinefun,split,sqrt,stars,start,stat,stem,step,stop,%
      storage,strstrheight,stripplot,strsplit,structure,strwidth,sub,%
      subset,substitute,substr,substring,sum,summary,sunflowerplot,svd,%
      sweep,switch,symbol,symbols,symnum,sys,status,system,t,table,%
      tabulate,tan,tanh,tapply,tempfile,terms,terrain,tetragamma,text,%
      time,title,topo,trace,traceback,transform,tri,trigamma,trunc,try,%
      ts,tsp,typeof,unclass,undebug,undoc,union,unique,uniroot,unix,%
      unlink,unlist,unname,untrace,update,upper,url,UseMethod,var,%
      variable,vector,Version,vi,warning,warnings,weighted,weights,%
      which,while,window,write,\%x\%,x11,X11,xedit,xemacs,xinch,xor,%
      xpdrows,xy,xyinch,yinch,zapsmall,zip},%
   otherkeywords={!,!=,~,$,*,\%,\&,\%/\%,\%*\%,\%\%,<-,<<-,_,/},%
   alsoother={._$},%
   sensitive,%
   morecomment=[l]\#,%
   morestring=[d]",%
   morestring=[d]'% 2001 Robert Denham
  }%

% Custom color for listings
\definecolor{codegreen}{rgb}{0,0.6,0}
\definecolor{codegray}{rgb}{0.5,0.5,0.5}
\definecolor{codepurple}{rgb}{0.58,0,0.82}
\definecolor{backcolour}{rgb}{0.95,0.95,0.92}

%------------------------------------------------------------
% Listings configuration for Julia
%------------------------------------------------------------
% Define the Julia language for the listings package
\lstdefinelanguage{Julia}
{
  morekeywords={abstract,baremodule,begin,break,catch,const,continue,do,else,elseif,end,export,false,finally,for,function,global,if,import,let,local,macro,module,quote,return,struct,true,try,using,while,in},
  sensitive=true,
  comment=[l]{\#},
  morecomment=[n]{\#=}{=\#}, % For block comments like #= ... =#
  string=[b]",
  morestring=[b]',
  alsoletter={!}, % Allow ! in identifiers for mutating functions
}

% Define the style for the code blocks
\lstdefinestyle{JuliaStyle}{
    backgroundcolor=\color{backcolour},   
    commentstyle=\color{codegreen},
    keywordstyle=\color{magenta},
    numberstyle=\tiny\color{codegray},
    stringstyle=\color{codepurple},
    basicstyle=\ttfamily\footnotesize,
    breakatwhitespace=false,         
    breaklines=true,                 
    captionpos=b,                    
    keepspaces=true,                 
    % numbers=left,                    
    % numbersep=5pt,                  
    showspaces=false,                
    showstringspaces=false,
    showtabs=false,                  
    tabsize=2,
    language=Julia % Use the newly defined Julia language
}
\lstset{style=JuliaStyle}

%%%%%%% 
%% Definitions in yellow boxes
\usepackage{etoolbox}
\setbeamercolor{block title}{use=structure,fg=structure.fg,bg=structure.fg!40!bg}
\setbeamercolor{block body}{parent=normal text,use=block title,bg=block title.bg!20!bg}

% Important definition
\BeforeBeginEnvironment{importantdefinition}{%
	\setbeamercolor{block title}{fg=black,bg=yellow!20!white}
	\setbeamercolor{block body}{fg=black, bg=yellow!05!white}
}
\AfterEndEnvironment{importantdefinition}{
	\setbeamercolor{block title}{use=structure,fg=structure.fg,bg=structure.fg!20!bg}
	\setbeamercolor{block body}{parent=normal text,use=block title,bg=block title.bg!50!bg, fg=black}
}
% Regular (less important) definition
\BeforeBeginEnvironment{definition}{%
	\setbeamercolor{block title}{fg=black,bg=black!15!white}
	\setbeamercolor{block body}{fg=black, bg=black!05!white}
}
\AfterEndEnvironment{definition}{
	\setbeamercolor{block title}{use=structure,fg=structure.fg,bg=structure.fg!20!bg}
	\setbeamercolor{block body}{parent=normal text,use=block title,bg=block title.bg!50!bg, fg=black}
}
% Important theorem (red)
\BeforeBeginEnvironment{importanttheorem}{%
	\setbeamercolor{block title}{fg=black,bg=red!20!white}
	\setbeamercolor{block body}{fg=black, bg=red!05!white}
}
\AfterEndEnvironment{importanttheorem}{
	\setbeamercolor{block title}{use=structure,fg=structure.fg,bg=structure.fg!20!bg}
	\setbeamercolor{block body}{parent=normal text,use=block title,bg=block title.bg!50!bg, fg=black}
}
\BeforeBeginEnvironment{importantproperty}{%
	\setbeamercolor{block title}{fg=black,bg=red!50!white}
	\setbeamercolor{block body}{fg=black, bg=red!30!white}
}
\AfterEndEnvironment{importantproperty}{
	\setbeamercolor{block title}{use=structure,fg=structure.fg,bg=structure.fg!20!bg}
	\setbeamercolor{block body}{parent=normal text,use=block title,bg=block title.bg!50!bg, fg=black}
}

%% Colours for title page, outline and the like
\definecolor{UMbrown}{RGB}{79,44,29}
\definecolor{HighlightGold}{RGB}{230,165,83}
\definecolor{DeepGreen}{RGB}{27,46,28}
\definecolor{MediumGreen}{RGB}{52,128,56}
\definecolor{LightGreen}{RGB}{150,224,154}
\definecolor{AccentRed}{RGB}{196,44,27}
%% These are the titles used in practice
% Title page
\colorlet{title_page_title_colour}{UMbrown}
\colorlet{title_page_author_colour}{UMbrown}
\colorlet{title_page_date_colour}{UMbrown}
\colorlet{title_page_institute_colour}{UMbrown}
% Outline page
\colorlet{outline_page_title_colour}{DeepGreen}
\colorlet{outline_page_list_colour}{UMbrown}
% Section page
\colorlet{section_page_title_colour}{DeepGreen}
\colorlet{section_page_list_colour}{UMbrown}
% Subection page
\colorlet{subsection_page_title_colour}{DeepGreen}
\colorlet{subsection_page_list_colour}{UMbrown}



\newtheorem{importantproposition}[theorem]{Proposition}
\BeforeBeginEnvironment{importantproposition}{%
	\setbeamercolor{block title}{fg=black,bg=red!20!white}
	\setbeamercolor{block body}{fg=black, bg=red!05!white}
}
\AfterEndEnvironment{importantproposition}{
	\setbeamercolor{block title}{use=structure,fg=structure.fg,bg=structure.fg!20!bg}
	\setbeamercolor{block body}{parent=normal text,use=block title,bg=block title.bg!50!bg, fg=black}
}
\IfFileExists{upquote.sty}{\usepackage{upquote}}{}
\begin{document}

% Set up cross-references and counter persistence
\setcounter{theorem}{154}

%%%%%%%%%%%%%%%%%%%%%%%%%%%%%%%%%
%%%%%%%%%%%%%%%%%%%%%%%%%%%%%%%%%
%% TITLE AND OUTLINE
%%%%%%%%%%%%%%%%%%%%%%%%%%%%%%%%%
%%%%%%%%%%%%%%%%%%%%%%%%%%%%%%%%%
\SStitlepage
\SSoutlinepage


\begin{frame}{From directed rankings to undirected communities}
In Units 08 and 09, we studied directed graphs, non-negative matrices and the Perron-Frobenius theorem
\vfill
There, the dominant eigenvector $\bv_1 > \b0$ associated with $\rho(A)$ answered a \emph{ranking} question: which vertices are most central or visited most often?
\vfill
In this unit, we turn to \defword{undirected graphs} and \defword{similarity networks}, where the associated matrices are \emph{real and symmetric}
\vfill
The main question shifts from ranking individual vertices to \defword{clustering}: how do we partition a graph (or a dataset) into cohesive communities with minimal connections between them?
\end{frame}


\begin{frame}{Why real symmetric matrices change everything}
When a matrix $M \in \M_n(\IR)$ is symmetric ($M^T = M$), the spectral theorem gives us three powerful tools that general directed graphs lack:
\vfill
\begin{enumerate}
\item Every eigenvalue $\lambda_1, \ldots, \lambda_n$ of $M$ is \textbf{real}
\vfill
\item Eigenvectors from distinct eigenspaces are \textbf{mutually orthogonal}, and $\IR^n$ has an \textbf{orthonormal basis of eigenvectors} $\{\bv_1, \ldots, \bv_n\}$
\vfill
\item Eigenvalues and eigenvectors are characterised variationally through the \textbf{quadratic form} $\bx^T M \bx$ and the \textbf{Rayleigh quotient} $\dfrac{\bx^T M \bx}{\bx^T \bx}$
\end{enumerate}
\end{frame}


\begin{frame}{Roadmap of Unit 10}
\begin{enumerate}
\item \textbf{Undirected graphs and similarity graphs:} turning point clouds $\bx_1, \ldots, \bx_n \in \IR^d$ into symmetric adjacency and weight matrices
\vfill
\item \textbf{The combinatorial graph Laplacian $L = D - A$:} degree matrix, incidence factorisation $L = B^T B$ and the quadratic form $\bx^T L \bx = \sum_{\{i,j\}\in E} a_{ij}(x_i - x_j)^2$
\vfill
\item \textbf{Spectral properties of $L$:} positive semi-definiteness, $\lambda_1 = 0$ with eigenvector $\11$ and why the multiplicity of $\lambda = 0$ counts connected components
\vfill
\item \textbf{Graph cuts, the Fiedler vector and spectral clustering:} expressing $\mathsf{cut}(S, \bar{S}) = \frac{1}{4}\bx^T L \bx$, spectral relaxation and partitioning with the Fiedler vector $\bv_2$
\end{enumerate}
\end{frame}


%%%%%%%%%%%%%%%%%%%%%
%%%%%%%%%%%%%%%%%%%%%
\SSsection{Undirected graphs and similarity graphs}

\begin{frame}{Undirected simple graphs}
Let $G = (V, E)$ be an undirected simple graph with vertex set $V = \{1, 2, \ldots, n\}$ and edge set $E$ consisting of unordered pairs $\{i, j\}$ of distinct vertices
\vfill
\begin{itemize}
\item Undirected means $\{i, j\} = \{j, i\}$: if $i$ is connected to $j$, then $j$ is connected to $i$
\vfill
\item Simple means there are no self-loops ($\{i, i\} \notin E$) and at most one edge between any pair of vertices
\end{itemize}
\vfill
\begin{importantdefinition}[Adjacency matrix of an undirected graph]
The \defword{adjacency matrix} of an undirected simple graph $G = (V, E)$ on $n$ vertices is the matrix $A = [a_{ij}] \in \M_n(\{0, 1\})$ defined by
\[
a_{ij} = \begin{cases}
1 & \text{if } \{i, j\} \in E \\
0 & \text{otherwise}
\end{cases}
\]
\end{importantdefinition}
\end{frame}


\begin{frame}{Symmetry of the adjacency matrix}
Because $\{i, j\} \in E \iff \{j, i\} \in E$ and $\{i, i\} \notin E$, the adjacency matrix $A$ of an undirected simple graph always satisfies
\[
A^T = A \quad (a_{ij} = a_{ji} \;\;\forall i, j) \qquad \text{and} \qquad a_{ii} = 0 \quad (\forall i = 1, \ldots, n)
\]
\vfill
\begin{columns}[c]
\begin{column}{0.45\textwidth}
\centering
\begin{tikzpicture}[scale=1.1, every node/.style={circle, draw=black, thick, fill=blue!10, minimum size=7mm, font=\small\bfseries}]
\node (1) at (0, 1.4) {1};
\node (2) at (0, 0) {2};
\node (3) at (1.6, 0.7) {3};
\node (4) at (3.2, 0.7) {4};
\draw[very thick] (1) -- (2);
\draw[very thick] (1) -- (3);
\draw[very thick] (2) -- (3);
\draw[very thick] (3) -- (4);
\end{tikzpicture}
\end{column}
\begin{column}{0.55\textwidth}
\[
A = \begin{pmatrix}
0 & 1 & 1 & 0 \\
1 & 0 & 1 & 0 \\
1 & 1 & 0 & 1 \\
0 & 0 & 1 & 0
\end{pmatrix} = A^T
\]
\end{column}
\end{columns}
\vfill
Notice that the trace of $A$ is $\tr(A) = \sum_{i=1}^n a_{ii} = 0$, so the $n$ real eigenvalues of $A$ always sum to $0$ (some are positive and some are negative)
\end{frame}


\begin{frame}{Weighted similarity graphs}
In data science, edges often carry a nonnegative weight measuring the \defword{similarity} (or affinity) between vertices $i$ and $j$
\vfill
\begin{definition}[Weighted similarity graph and weight matrix]
A \defword{weighted undirected graph} $G = (V, E, W)$ has a symmetric weight matrix $W = [w_{ij}] \in \M_n(\IRplus)$ satisfying
\[
w_{ij} = w_{ji} \ge 0 \quad (\forall i, j = 1, \ldots, n) \qquad \text{and} \qquad w_{ii} = 0 \quad (\forall i = 1, \ldots, n)
\]
Here $w_{ij} > 0$ if $\{i, j\} \in E$ and $w_{ij} = 0$ if $\{i, j\} \notin E$
\end{definition}
\vfill
\begin{itemize}
\item Large $w_{ij}$ means vertices $i$ and $j$ are \textbf{strongly connected} (highly similar)
\item $w_{ij} = 0$ (or $w_{ij} \approx 0$) means vertices $i$ and $j$ are \textbf{unconnected} (dissimilar)
\item An unweighted graph is simply the special case $W = A \in \M_n(\{0, 1\})$
\end{itemize}
\end{frame}


\begin{frame}{Why turn a dataset into a similarity graph?}
Suppose we are given $n$ data points $\bx_1, \bx_2, \ldots, \bx_n \in \IR^d$ and want to partition them into clusters
\vfill
\begin{columns}[c]
\begin{column}{0.56\textwidth}
\begin{itemize}
\item Centroid-based methods ($k$-means) and linear projections (PCA) rely on \textbf{global Euclidean distances} $\|\bx_i - \bx_j\|_2$ and linear separating hyperplanes
\vfill
\item They work well when clusters are compact, convex clouds
\vfill
\item They fail completely when clusters have \textbf{non-convex geometry} (such as concentric rings, spirals or interlocking crescents), where points in the same cluster are far apart in $\IR^d$ but connected through a chain of nearby neighbours
\end{itemize}
\end{column}
\begin{column}{0.44\textwidth}
\centering
\begin{tikzpicture}[scale=0.95]
% Outer ring
\foreach \ang in {0,30,...,330} {
  \filldraw[blue!75!black] (\ang:1.6) circle (2pt);
}
\draw[blue!40, thick, dashed] (0,0) circle (1.6);
% Inner ring
\foreach \ang in {15,75,...,315} {
  \filldraw[red!80!black] (\ang:0.65) circle (2pt);
}
\draw[red!40, thick, dashed] (0,0) circle (0.65);
% Linear cut line
\draw[thick, black!70, dashdotted] (-1.9,-0.8) -- (1.9,0.8) node[right, font=\scriptsize] {linear cut};
\end{tikzpicture}
\end{column}
\end{columns}
\vfill
Representing the data as a \textbf{local similarity graph} turns non-convex geometric clustering into a \textbf{graph connectivity} problem
\end{frame}


\begin{frame}{Constructing similarity graphs from data points}
Given data points $\bx_1, \ldots, \bx_n \in \IR^d$ with pairwise distances $d(\bx_i, \bx_j) = \|\bx_i - \bx_j\|_2$, there are three standard ways to build a symmetric similarity graph:
\vfill
\begin{enumerate}
\item \defword{$\varepsilon$-neighbourhood graph}: Fix a threshold $\varepsilon > 0$ and connect $i \neq j$ with an unweighted edge ($a_{ij} = 1$) if and only if
\[
\|\bx_i - \bx_j\|_2 \le \varepsilon
\]
\item \defword{$k$-nearest neighbour ($k$-NN) graph}: Connect $i$ and $j$ if $\bx_i$ is among the $k$ closest points to $\bx_j$ \emph{or} $\bx_j$ is among the $k$ closest points to $\bx_i$ (symmetrised so $A^T = A$)
\vfill
\item \defword{Gaussian (RBF) similarity graph}: Connect all pairs $i \neq j$ with weight
\[
w_{ij} = \exp\!\left(-\frac{\|\bx_i - \bx_j\|_2^2}{2\sigma^2}\right), \qquad w_{ii} = 0
\]
where $\sigma > 0$ is the bandwidth (scale) parameter
\end{enumerate}
\end{frame}


\begin{frame}{Example: from 4 points in $\IR^2$ to a similarity graph}
Consider $n = 4$ points in $\IR^2$: two nearby points $\bx_1 = (0, 0)^T, \bx_2 = (0, 1)^T$ on the left and two nearby points $\bx_3 = (3, 0)^T, \bx_4 = (3, 1)^T$ on the right
\vfill
\begin{columns}[c]
\begin{column}{0.48\textwidth}
Pairwise Euclidean distances $\|\bx_i - \bx_j\|_2$:
\[
\Delta = \begin{pmatrix}
0 & 1 & 3 & \sqrt{10} \\
1 & 0 & \sqrt{10} & 3 \\
3 & \sqrt{10} & 0 & 1 \\
\sqrt{10} & 3 & 1 & 0
\end{pmatrix}
\]
With $\varepsilon = 1.5$, the \textbf{$\varepsilon$-neighbourhood graph} has adjacency matrix
\[
A = \begin{pmatrix}
0 & 1 & 0 & 0 \\
1 & 0 & 0 & 0 \\
0 & 0 & 0 & 1 \\
0 & 0 & 1 & 0
\end{pmatrix}
\]
\end{column}
\begin{column}{0.52\textwidth}
With $\sigma = 1$, the \textbf{Gaussian similarity matrix} $w_{ij} = e^{-\|\bx_i - \bx_j\|_2^2 / 2}$ ($i \neq j$) is
\[
W \approx \begin{pmatrix}
0 & 0.607 & 0.011 & 0.007 \\
0.607 & 0 & 0.007 & 0.011 \\
0.011 & 0.007 & 0 & 0.607 \\
0.007 & 0.011 & 0.607 & 0
\end{pmatrix}
\]
\begin{itemize}
\item Within-cluster weights: $e^{-1/2} \approx 0.607$
\item Between-cluster weights: $e^{-9/2} \approx 0.011$ and $e^{-10/2} \approx 0.007$
\end{itemize}
\end{column}
\end{columns}
\end{frame}


\begin{frame}{Role of the scale parameter $\varepsilon$, $k$ or $\sigma$}
In all three constructions, the scale parameter controls which pairs are treated as local neighbours:
\vfill
\begin{columns}[c]
\begin{column}{0.32\textwidth}
\centering
\textbf{Scale too small}\\[4pt]
\begin{tikzpicture}[scale=0.8]
\node[circle,fill=blue!70,inner sep=2pt] (a1) at (0,1) {};
\node[circle,fill=blue!70,inner sep=2pt] (a2) at (0.6,1.2) {};
\node[circle,fill=blue!70,inner sep=2pt] (a3) at (0.3,0.4) {};
\node[circle,fill=red!70,inner sep=2pt] (b1) at (2.2,1) {};
\node[circle,fill=red!70,inner sep=2pt] (b2) at (2.8,1.2) {};
\node[circle,fill=red!70,inner sep=2pt] (b3) at (2.5,0.4) {};
\draw[thick] (a1) -- (a2);
\draw[thick] (b1) -- (b2);
\end{tikzpicture}\\[4pt]
\footnotesize True clusters shatter into many isolated pieces
\end{column}
\begin{column}{0.34\textwidth}
\centering
\textbf{Well-chosen scale}\\[4pt]
\begin{tikzpicture}[scale=0.8]
\node[circle,fill=blue!70,inner sep=2pt] (a1) at (0,1) {};
\node[circle,fill=blue!70,inner sep=2pt] (a2) at (0.6,1.2) {};
\node[circle,fill=blue!70,inner sep=2pt] (a3) at (0.3,0.4) {};
\node[circle,fill=red!70,inner sep=2pt] (b1) at (2.2,1) {};
\node[circle,fill=red!70,inner sep=2pt] (b2) at (2.8,1.2) {};
\node[circle,fill=red!70,inner sep=2pt] (b3) at (2.5,0.4) {};
\draw[thick] (a1) -- (a2) -- (a3) -- (a1);
\draw[thick] (b1) -- (b2) -- (b3) -- (b1);
\draw[dashed, gray] (a2) -- (b1);
\end{tikzpicture}\\[4pt]
\footnotesize Dense within clusters, sparse or zero between clusters
\end{column}
\begin{column}{0.32\textwidth}
\centering
\textbf{Scale too large}\\[4pt]
\begin{tikzpicture}[scale=0.8]
\node[circle,fill=blue!70,inner sep=2pt] (a1) at (0,1) {};
\node[circle,fill=blue!70,inner sep=2pt] (a2) at (0.6,1.2) {};
\node[circle,fill=blue!70,inner sep=2pt] (a3) at (0.3,0.4) {};
\node[circle,fill=red!70,inner sep=2pt] (b1) at (2.2,1) {};
\node[circle,fill=red!70,inner sep=2pt] (b2) at (2.8,1.2) {};
\node[circle,fill=red!70,inner sep=2pt] (b3) at (2.5,0.4) {};
\foreach \u in {a1,a2,a3}
  \foreach \v in {b1,b2,b3}
    \draw[thin, gray!70] (\u) -- (\v);
\draw[thick] (a1) -- (a2) -- (a3) -- (a1);
\draw[thick] (b1) -- (b2) -- (b3) -- (b1);
\end{tikzpicture}\\[4pt]
\footnotesize Shortcut edges merge distinct clusters together
\end{column}
\end{columns}
\vfill
In what follows, we write $A = [a_{ij}]$ for the adjacency or similarity matrix (whether $a_{ij} \in \{0, 1\}$ or $a_{ij} = w_{ij} \ge 0$): all formulas hold identically in both cases
\end{frame}


\begin{frame}[fragile]{R illustration: two concentric rings}
Let us generate $n = 160$ points on two concentric circles in $\IR^2$ and compare $k$-means in $\IR^2$ against an $\varepsilon$-neighbourhood graph
\vfill
\begin{knitrout}
\definecolor{shadecolor}{rgb}{0.969, 0.969, 0.969}\color{fgcolor}\begin{kframe}
\begin{alltt}
\hlkwd{set.seed}\hldef{(}\hlnum{2740}\hldef{)}
\hldef{n1} \hlkwb{=} \hlnum{80}\hldef{; n2} \hlkwb{=} \hlnum{80}
\hldef{theta1} \hlkwb{=} \hlkwd{seq}\hldef{(}\hlnum{0}\hldef{,} \hlnum{2} \hlopt{*} \hldef{pi,} \hlkwc{length.out} \hldef{= n1} \hlopt{+} \hlnum{1}\hldef{)[}\hlopt{-}\hldef{(n1} \hlopt{+} \hlnum{1}\hldef{)]}
\hldef{theta2} \hlkwb{=} \hlkwd{seq}\hldef{(}\hlnum{0}\hldef{,} \hlnum{2} \hlopt{*} \hldef{pi,} \hlkwc{length.out} \hldef{= n2} \hlopt{+} \hlnum{1}\hldef{)[}\hlopt{-}\hldef{(n2} \hlopt{+} \hlnum{1}\hldef{)]}
\hldef{X} \hlkwb{=} \hlkwd{rbind}\hldef{(}
  \hlkwd{cbind}\hldef{(}\hlnum{1.0} \hlopt{*} \hlkwd{cos}\hldef{(theta1),} \hlnum{1.0} \hlopt{*} \hlkwd{sin}\hldef{(theta1))} \hlopt{+} \hlkwd{rnorm}\hldef{(}\hlnum{2} \hlopt{*} \hldef{n1,} \hlkwc{sd} \hldef{=} \hlnum{0.06}\hldef{),}
  \hlkwd{cbind}\hldef{(}\hlnum{2.3} \hlopt{*} \hlkwd{cos}\hldef{(theta2),} \hlnum{2.3} \hlopt{*} \hlkwd{sin}\hldef{(theta2))} \hlopt{+} \hlkwd{rnorm}\hldef{(}\hlnum{2} \hlopt{*} \hldef{n2,} \hlkwc{sd} \hldef{=} \hlnum{0.06}\hldef{)}
\hldef{)}
\hlcom{# epsilon-neighbourhood adjacency matrix (epsilon = 0.65)}
\hldef{dist_mat} \hlkwb{=} \hlkwd{as.matrix}\hldef{(}\hlkwd{dist}\hldef{(X))}
\hldef{A_rings} \hlkwb{=} \hldef{(dist_mat} \hlopt{<=} \hlnum{0.65} \hlopt{&} \hldef{dist_mat} \hlopt{>} \hlnum{0}\hldef{)} \hlopt{*} \hlnum{1}
\end{alltt}
\end{kframe}
\end{knitrout}
\end{frame}




\maxFrameImage{FIGS/U10-plot-concentric-rings-1.pdf}


%%%%%%%%%%%%%%%%%%%%%
%%%%%%%%%%%%%%%%%%%%%
\SSsection{The combinatorial graph Laplacian}

\SSsubsection{Definition and basic examples}

\begin{frame}{Vertex degrees and the degree matrix}
\begin{importantdefinition}[Vertex degree and degree matrix]
Let $G = (V, E)$ be an undirected graph on $n$ vertices with symmetric adjacency (or weight) matrix $A = [a_{ij}] \in \M_n(\IRplus)$. The \defword{degree} of vertex $i$ is the row sum
\[
d_i = \sum_{j=1}^n a_{ij} = (A\11)_i \qquad (i = 1, \ldots, n)
\]
where $\11 = (1, 1, \ldots, 1)^T \in \IR^n$. The \defword{degree matrix} of $G$ is the diagonal matrix
\[
D = \diag(d_1, d_2, \ldots, d_n) = \begin{pmatrix}
d_1 & 0 & \cdots & 0 \\
0 & d_2 & \cdots & 0 \\
\vdots & \vdots & \ddots & \vdots \\
0 & 0 & \cdots & d_n
\end{pmatrix}
\]
\end{importantdefinition}
\end{frame}


\begin{frame}{The combinatorial graph Laplacian}
\begin{importantdefinition}[Combinatorial graph Laplacian]
Let $G = (V, E)$ be an undirected graph on $n$ vertices with degree matrix $D$ and adjacency (or weight) matrix $A$. The \defword{combinatorial graph Laplacian} (or simply \defword{Laplacian matrix}) of $G$ is the $n \times n$ matrix
\[
L = D - A
\]
Entry-by-entry, $L = [\ell_{ij}]$ is given by
\[
\ell_{ij} = \begin{cases}
d_i & \text{if } i = j \\
-a_{ij} & \text{if } i \neq j
\end{cases}
\]
\end{importantdefinition}
\vfill
For a simple unweighted graph, $\ell_{ii} = d_i$, $\ell_{ij} = -1$ if $\{i, j\} \in E$ and $\ell_{ij} = 0$ otherwise
\end{frame}


\begin{frame}{Example 1: Laplacian of a 4-vertex paw graph}
Consider the 4-vertex graph $G = (V, E)$ below (a triangle $\{1, 2, 3\}$ with a pendant edge $\{3, 4\}$):
\vfill
\begin{columns}[c]
\begin{column}{0.36\textwidth}
\centering
\begin{tikzpicture}[scale=1.15, every node/.style={circle, draw=black, thick, fill=yellow!20, minimum size=7mm, font=\small\bfseries}]
\node (1) at (0, 1.4) {1};
\node (2) at (0, 0) {2};
\node (3) at (1.4, 0.7) {3};
\node (4) at (2.8, 0.7) {4};
\draw[very thick] (1) -- (2);
\draw[very thick] (1) -- (3);
\draw[very thick] (2) -- (3);
\draw[very thick] (3) -- (4);
\end{tikzpicture}
\vskip0.2cm
\footnotesize
Degrees: $d_1 = 2$, $d_2 = 2$, $d_3 = 3$, $d_4 = 1$
\end{column}
\begin{column}{0.64\textwidth}
\[
D = \begin{pmatrix}
2 & 0 & 0 & 0 \\
0 & 2 & 0 & 0 \\
0 & 0 & 3 & 0 \\
0 & 0 & 0 & 1
\end{pmatrix}, \qquad
A = \begin{pmatrix}
0 & 1 & 1 & 0 \\
1 & 0 & 1 & 0 \\
1 & 1 & 0 & 1 \\
0 & 0 & 1 & 0
\end{pmatrix}
\]
\[
L = D - A = \begin{pmatrix}
2 & -1 & -1 & 0 \\
-1 & 2 & -1 & 0 \\
-1 & -1 & 3 & -1 \\
0 & 0 & -1 & 1
\end{pmatrix}
\]
\end{column}
\end{columns}
\vfill
Quick sanity check on $L$:
\begin{itemize}
\item $L$ is symmetric ($L^T = L$), with positive degrees on the diagonal and $-1$ on edges
\item Every row and every column of $L$ sums to $0$: e.g. row 3 has $(-1) + (-1) + 3 + (-1) = 0$
\end{itemize}
\end{frame}


\begin{frame}{Example 2: Laplacian of a 5-vertex bowtie graph}
Consider the 5-vertex bowtie graph formed by two triangles $\{1, 2, 3\}$ and $\{3, 4, 5\}$ sharing central vertex $3$:
\vfill
\begin{columns}[c]
\begin{column}{0.38\textwidth}
\centering
\begin{tikzpicture}[scale=1.05, every node/.style={circle, draw=black, thick, fill=yellow!20, minimum size=7mm, font=\small\bfseries}]
\node (1) at (-1.5, 0.8) {1};
\node (2) at (-1.5, -0.8) {2};
\node (3) at (0, 0) {3};
\node (4) at (1.5, 0.8) {4};
\node (5) at (1.5, -0.8) {5};
\draw[very thick] (1) -- (2) -- (3) -- (1);
\draw[very thick] (3) -- (4) -- (5) -- (3);
\end{tikzpicture}
\vskip0.2cm
\footnotesize
Degrees: $d_1 = d_2 = d_4 = d_5 = 2$, $d_3 = 4$
\end{column}
\begin{column}{0.62\textwidth}
\[
L = D - A = \begin{pmatrix}
2 & -1 & -1 & 0 & 0 \\
-1 & 2 & -1 & 0 & 0 \\
-1 & -1 & 4 & -1 & -1 \\
0 & 0 & -1 & 2 & -1 \\
0 & 0 & -1 & -1 & 2
\end{pmatrix}
\]
\end{column}
\end{columns}
\vfill
Notice how the structure of $L$ reflects the graph topology:
\begin{itemize}
\item The upper-left $2\times 2$ block $\{1, 2\}$ and lower-right $2\times 2$ block $\{4, 5\}$ have no direct edges between them ($0$ in the off-diagonal blocks)
\item All paths between $\{1, 2\}$ and $\{4, 5\}$ pass through vertex $3$ (row and column $3$)
\end{itemize}
\end{frame}


\begin{frame}{Example 3: Laplacian of a weighted similarity graph}
For a weighted graph $G = (V, E, W)$, the exact same recipe $L = D - W$ applies using weighted degrees $d_i = \sum_{j=1}^n w_{ij}$:
\vfill
\begin{columns}[c]
\begin{column}{0.42\textwidth}
\centering
\begin{tikzpicture}[scale=1.15]
\node[circle, draw=black, thick, fill=yellow!20, minimum size=7mm, font=\small\bfseries] (1) at (0, 1.4) {1};
\node[circle, draw=black, thick, fill=yellow!20, minimum size=7mm, font=\small\bfseries] (2) at (0, 0) {2};
\node[circle, draw=black, thick, fill=yellow!20, minimum size=7mm, font=\small\bfseries] (3) at (1.6, 0.7) {3};
\node[circle, draw=black, thick, fill=yellow!20, minimum size=7mm, font=\small\bfseries] (4) at (3.2, 0.7) {4};
\draw[very thick] (1) -- node[left, font=\footnotesize, text=blue!80!black] {$5$} (2);
\draw[very thick] (1) -- node[above right, font=\footnotesize, text=blue!80!black] {$2$} (3);
\draw[very thick] (2) -- node[below right, font=\footnotesize, text=blue!80!black] {$4$} (3);
\draw[very thick] (3) -- node[above, font=\footnotesize, text=red!80!black] {$0.5$} (4);
\end{tikzpicture}
\vskip0.2cm
\footnotesize
Weighted degrees:\\
$d_1 = 5 + 2 = 7$, $\; d_2 = 5 + 4 = 9$\\
$d_3 = 2 + 4 + 0.5 = 6.5$, $\; d_4 = 0.5$
\end{column}
\begin{column}{0.58\textwidth}
\[
W = \begin{pmatrix}
0 & 5 & 2 & 0 \\
5 & 0 & 4 & 0 \\
2 & 4 & 0 & 0.5 \\
0 & 0 & 0.5 & 0
\end{pmatrix}
\]
\[
L = D - W = \begin{pmatrix}
7 & -5 & -2 & 0 \\
-5 & 9 & -4 & 0 \\
-2 & -4 & 6.5 & -0.5 \\
0 & 0 & -0.5 & 0.5
\end{pmatrix}
\]
\end{column}
\end{columns}
\vfill
Every row and column of $L$ still sums to $0$, and $\tr(L) = \sum_{i=1}^n d_i = 2\sum_{\{i,j\}\in E} w_{ij} = 23$
\end{frame}


\begin{frame}{Connection with the oriented incidence matrix}
We can also factor $L$ as a Gram matrix $B^T B$ by assigning an arbitrary orientation to each edge of $G = (V, E)$
\vfill
\begin{definition}[Oriented edge-vertex incidence matrix]
Let $G = (V, E)$ have $n = |V|$ vertices and $m = |E|$ edges $e_1, \ldots, e_m$. Choose an arbitrary direction $e_k = (i, j)$ for each edge ($i \to j$). The \defword{oriented incidence matrix} $B = [b_{kv}] \in \M_{m, n}(\{-1, 0, 1\})$ has one row per edge and one column per vertex:
\[
b_{kv} = \begin{cases}
+1 & \text{if } v = i \text{ (tail of } e_k\text{)} \\
-1 & \text{if } v = j \text{ (head of } e_k\text{)} \\
0 & \text{otherwise}
\end{cases}
\]
\end{definition}
\vfill
\begin{proposition}[Incidence factorisation of the Laplacian]
\label{prop:incidence-laplacian}
For any choice of edge orientations on an unweighted graph $G$, $L = B^T B$
\end{proposition}
\end{frame}


\begin{frame}[red]{Proof of Proposition~\ref{prop:incidence-laplacian} ($L = B^T B$) and worked check}
\textbf{Proof:} The $(u, v)$-entry of $B^T B \in \M_n(\IR)$ is the dot product of columns $u$ and $v$ of $B$:
\[
(B^T B)_{uv} = \sum_{k=1}^m b_{ku} b_{kv}
\]
\begin{itemize}
\item If $u = v$: $b_{ku}^2 = (\pm 1)^2 = 1$ whenever edge $e_k$ is incident to vertex $u$, and $0$ otherwise, so $(B^T B)_{uu} = \sum_{k=1}^m b_{ku}^2 = d_u = \ell_{uu}$
\item If $u \neq v$: the product $b_{ku} b_{kv}$ is non-zero only if edge $e_k$ connects $u$ and $v$, in which case one entry is $+1$ and the other is $-1$, so $b_{ku} b_{kv} = -1 = \ell_{uv}$ \hfill $\qed$
\end{itemize}
\vfill
\textbf{Check on the 4-vertex paw graph} (orienting $e_1=(1,2), e_2=(1,3), e_3=(2,3), e_4=(3,4)$):
\[
B = \begin{pmatrix}
1 & -1 & 0 & 0 \\
1 & 0 & -1 & 0 \\
0 & 1 & -1 & 0 \\
0 & 0 & 1 & -1
\end{pmatrix}
\implies
B^T B = \begin{pmatrix}
2 & -1 & -1 & 0 \\
-1 & 2 & -1 & 0 \\
-1 & -1 & 3 & -1 \\
0 & 0 & -1 & 1
\end{pmatrix} = L
\]
\end{frame}


\begin{frame}{Elementary matrix properties of $L$}
From $L = D - A$ and $L = B^T B$, several useful matrix identities follow immediately:
\vfill
\begin{enumerate}
\item \textbf{Trace of $L$:} By the handshaking lemma, the sum of the eigenvalues of $L$ equals twice the total edge weight:
\[
\sum_{k=1}^n \lambda_k = \tr(L) = \sum_{i=1}^n d_i = 2 \sum_{\{i, j\} \in E} a_{ij} \qquad (= 2|E| \text{ for an unweighted graph})
\]
\item \textbf{Sum of rank-one edge matrices:} Each edge $e = \{i, j\} \in E$ contributes a $2\times 2$ sub-Laplacian on vertices $i$ and $j$. Letting $\be_1, \ldots, \be_n$ denote the standard basis of $\IR^n$,
\[
L = \sum_{\{i, j\} \in E} a_{ij}(\be_i - \be_j)(\be_i - \be_j)^T
\]
\item \textbf{Weak diagonal dominance:} Every row $i$ satisfies $|\ell_{ii}| = d_i = \sum_{j \neq i} a_{ij} = \sum_{j \neq i} |\ell_{ij}|$
\end{enumerate}
\end{frame}


%%%%%%%%%%%%%%%%%%%%%
\SSsubsection{The Laplacian quadratic form}

\begin{frame}{Vectors on $\IR^n$ as functions on the vertices}
How does $L$ act on a vector $\bx = (x_1, x_2, \ldots, x_n)^T \in \IR^n$?
\vfill
\begin{columns}[c]
\begin{column}{0.52\textwidth}
Think of $\bx \in \IR^n$ as assigning a real number (a ``signal'' or ``potential'') $x_i \in \IR$ to each vertex $i \in V$:
\begin{itemize}
\item At vertex $i$, the $i$th entry of $L\bx$ is
\[
(L\bx)_i = d_i x_i - \sum_{j=1}^n a_{ij} x_j = \sum_{j=1}^n a_{ij}(x_i - x_j)
\]
\item Along each edge $e_k = (i, j)$, the $k$th entry of $B\bx$ is the \textbf{edge difference}
\[
(B\bx)_k = x_i - x_j
\]
\end{itemize}
\end{column}
\begin{column}{0.48\textwidth}
\centering
\begin{tikzpicture}[scale=1.15]
\node[circle, draw=black, thick, fill=blue!15, minimum size=8mm, font=\small\bfseries] (1) at (0, 1.5) {$x_1$};
\node[circle, draw=black, thick, fill=blue!15, minimum size=8mm, font=\small\bfseries] (2) at (0, 0) {$x_2$};
\node[circle, draw=black, thick, fill=blue!15, minimum size=8mm, font=\small\bfseries] (3) at (1.7, 0.75) {$x_3$};
\node[circle, draw=black, thick, fill=blue!15, minimum size=8mm, font=\small\bfseries] (4) at (3.4, 0.75) {$x_4$};
\draw[very thick] (1) -- node[left, font=\scriptsize, text=red!80!black] {$x_1 - x_2$} (2);
\draw[very thick] (1) -- node[above right, font=\scriptsize, text=red!80!black] {$x_1 - x_3$} (3);
\draw[very thick] (2) -- node[below right, font=\scriptsize, text=red!80!black] {$x_2 - x_3$} (3);
\draw[very thick] (3) -- node[above, font=\scriptsize, text=red!80!black] {$x_3 - x_4$} (4);
\end{tikzpicture}
\end{column}
\end{columns}
\vfill
Multiplying on the left by $\bx^T$ produces the fundamental identity of spectral graph theory
\end{frame}


\begin{frame}{The Laplacian quadratic form}
\begin{importanttheorem}[Laplacian quadratic form]
\label{th:laplacian-quadratic-form}
Let $G = (V, E)$ be an undirected graph on $n$ vertices with symmetric weight matrix $A = [a_{ij}]$ ($a_{ij} = a_{ji} \ge 0$) and combinatorial Laplacian $L = D - A$. Then for every vector $\bx = (x_1, \ldots, x_n)^T \in \IR^n$,
\[
\bx^T L \bx = \frac{1}{2} \sum_{i=1}^n \sum_{j=1}^n a_{ij}(x_i - x_j)^2 = \sum_{\{i, j\} \in E} a_{ij}(x_i - x_j)^2
\]
In particular, for an unweighted simple graph ($a_{ij} \in \{0, 1\}$),
\[
\bx^T L \bx = \sum_{\{i, j\} \in E} (x_i - x_j)^2
\]
\end{importanttheorem}
\end{frame}


\begin{frame}[red]{Proof of Theorem~\ref{th:laplacian-quadratic-form} (Method 1: direct algebraic expansion)}
Let $\bx = (x_1, \ldots, x_n)^T \in \IR^n$. Using $L = D - A$ and $d_i = \sum_{j=1}^n a_{ij}$, we expand
\begin{align*}
\bx^T L \bx &= \bx^T D \bx - \bx^T A \bx = \sum_{i=1}^n d_i x_i^2 - \sum_{i=1}^n \sum_{j=1}^n a_{ij} x_i x_j \\
&= \sum_{i=1}^n \left(\sum_{j=1}^n a_{ij}\right) x_i^2 - \sum_{i=1}^n \sum_{j=1}^n a_{ij} x_i x_j
\end{align*}
By symmetry $a_{ij} = a_{ji}$, relabelling dummy indices $i \leftrightarrow j$ gives $\sum_{i,j} a_{ij} x_i^2 = \sum_{i,j} a_{ji} x_j^2 = \sum_{i,j} a_{ij} x_j^2$, so we can split the first sum into two equal halves:
\begin{align*}
\bx^T L \bx &= \frac{1}{2}\sum_{i=1}^n \sum_{j=1}^n a_{ij} x_i^2 - \sum_{i=1}^n \sum_{j=1}^n a_{ij} x_i x_j + \frac{1}{2}\sum_{i=1}^n \sum_{j=1}^n a_{ij} x_j^2 \\
&= \frac{1}{2}\sum_{i=1}^n \sum_{j=1}^n a_{ij}\left(x_i^2 - 2x_i x_j + x_j^2\right) = \frac{1}{2}\sum_{i=1}^n \sum_{j=1}^n a_{ij}(x_i - x_j)^2
\end{align*}
Since each undirected edge $\{i, j\} \in E$ appears twice in the double sum (as $(i, j)$ and $(j, i)$), the factor $\frac{1}{2}$ cancels when summing over $\{i, j\} \in E$ \hfill $\qed$
\end{frame}


\begin{frame}[red]{Proof of Theorem~\ref{th:laplacian-quadratic-form} (Method 2: via the incidence matrix)}
For an unweighted graph $G = (V, E)$, Proposition~\ref{prop:incidence-laplacian} gives $L = B^T B$, where $B \in \M_{|E|, n}$ is the oriented incidence matrix
\vfill
For any $\bx \in \IR^n$, the $k$th entry of $B\bx \in \IR^{|E|}$ corresponding to oriented edge $e_k = (i, j)$ is
\[
(B\bx)_k = x_i - x_j
\]
\vfill
Therefore,
\[
\bx^T L \bx = \bx^T (B^T B) \bx = (B\bx)^T (B\bx) = \|B\bx\|_2^2 = \sum_{k=1}^{|E|} (B\bx)_k^2 = \sum_{\{i, j\} \in E} (x_i - x_j)^2
\]
\vfill
(For a weighted graph with diagonal edge-weight matrix $W_E = \diag(w_{e_1}, \ldots, w_{e_{|E|}})$, writing $L = B^T W_E B = (W_E^{1/2} B)^T (W_E^{1/2} B)$ gives $\bx^T L \bx = \|W_E^{1/2} B\bx\|_2^2 = \sum_{\{i,j\}\in E} w_{ij}(x_i - x_j)^2$) \hfill $\qed$
\end{frame}


\begin{frame}{Geometric interpretation: smoothness (Dirichlet energy) on a graph}
The quadratic form $\bx^T L \bx = \sum_{\{i,j\}\in E} a_{ij}(x_i - x_j)^2$ is called the \defword{Dirichlet energy} of the vector $\bx$ on the graph $G$:
\vfill
\begin{itemize}
\item Every edge $\{i, j\} \in E$ acts like a spring with stiffness $a_{ij}$ that penalises the squared difference $(x_i - x_j)^2$ between its endpoints
\vfill
\item \textbf{Smooth vector} ($\bx^T L \bx$ small): adjacent vertices have nearly equal values ($x_i \approx x_j$ whenever $a_{ij}$ is large)
\vfill
\item \textbf{Constant vector} ($\bx^T L \bx = 0$): $x_i = x_j$ across every edge of $G$
\vfill
\item \textbf{Rough / oscillating vector} ($\bx^T L \bx$ large): adjacent vertices have very different values (e.g. alternating signs $+1$ and $-1$ across edges)
\end{itemize}
\end{frame}


\begin{frame}{Example: evaluating $\bx^T L \bx$ on the 4-vertex paw graph}
Take the 4-vertex paw graph with edges $\{1, 2\}, \{1, 3\}, \{2, 3\}, \{3, 4\}$ and compare two vectors of equal norm:
\vfill
\begin{columns}[t]
\begin{column}{0.48\textwidth}
\textbf{Case 1 (smooth):} $\bx = (1, 1, 0, -1)^T$
\begin{align*}
\bx^T L \bx &= (x_1 - x_2)^2 + (x_1 - x_3)^2 \\
&\quad + (x_2 - x_3)^2 + (x_3 - x_4)^2 \\
&= (1 - 1)^2 + (1 - 0)^2 \\
&\quad + (1 - 0)^2 + (0 - (-1))^2 \\
&= 0 + 1 + 1 + 1 = 3
\end{align*}
\footnotesize
Adjacent vertices $\{1, 2\}$ share the same value $1$, and values change gradually along the graph
\end{column}
\begin{column}{0.48\textwidth}
\textbf{Case 2 (oscillating):} $\by = (1, -1, 0, 1)^T$
\begin{align*}
\by^T L \by &= (y_1 - y_2)^2 + (y_1 - y_3)^2 \\
&\quad + (y_2 - y_3)^2 + (y_3 - y_4)^2 \\
&= (1 - (-1))^2 + (1 - 0)^2 \\
&\quad + (-1 - 0)^2 + (0 - 1)^2 \\
&= 4 + 1 + 1 + 1 = 7
\end{align*}
\footnotesize
Adjacent vertices $1$ and $2$ jump from $+1$ to $-1$, producing a much larger energy
\end{column}
\end{columns}
\end{frame}


%%%%%%%%%%%%%%%%%%%%%
%%%%%%%%%%%%%%%%%%%%%
\SSsection{Spectral properties of the Laplacian}

\begin{frame}{Fundamental spectral properties of the Laplacian}
\begin{importanttheorem}[Fundamental spectral properties of $L$]
\label{th:laplacian-spectrum}
Let $G = (V, E)$ be an undirected graph on $n$ vertices with symmetric nonnegative weight matrix $A = [a_{ij}]$ and combinatorial Laplacian $L = D - A$. Then:
\begin{enumerate}
\item $L$ is real and symmetric ($L^T = L$), so it has $n$ real eigenvalues and an orthonormal basis of real eigenvectors $\{\bv_1, \bv_2, \ldots, \bv_n\}$
\item $L$ is \defword{positive semi-definite} ($\bx^T L \bx \ge 0$ for all $\bx \in \IR^n$), so all eigenvalues of $L$ are nonnegative and can be ordered in ascending order:
\[
0 \le \lambda_1 \le \lambda_2 \le \cdots \le \lambda_n
\]
\item The all-ones vector $\11 = (1, 1, \ldots, 1)^T$ satisfies $L\11 = \b0$, so $\lambda_1 = 0$ is always the smallest eigenvalue of $L$ with eigenvector $\bv_1 = \11$ (or unit eigenvector $\frac{1}{\sqrt{n}}\11$)
\end{enumerate}
\end{importanttheorem}
\end{frame}


\begin{frame}[red]{Proof of Theorem~\ref{th:laplacian-spectrum}}
\begin{enumerate}
\item \textbf{Symmetry and real spectrum:} Since $D$ is diagonal ($D^T = D$) and $A$ is symmetric ($A^T = A$),
\[
L^T = (D - A)^T = D^T - A^T = D - A = L
\]
By the spectral theorem for real symmetric matrices, $L$ has $n$ real eigenvalues and an orthonormal basis of real eigenvectors in $\IR^n$
\vfill
\item \textbf{Positive semi-definiteness and $\lambda_k \ge 0$:} By Theorem~\ref{th:laplacian-quadratic-form}, for any $\bx \in \IR^n$,
\[
\bx^T L \bx = \sum_{\{i, j\} \in E} a_{ij}(x_i - x_j)^2 \ge 0
\]
since $a_{ij} \ge 0$ and $(x_i - x_j)^2 \ge 0$. If $L\bv_k = \lambda_k \bv_k$ with $\bv_k \neq \b0$, then
\[
\lambda_k = \frac{\bv_k^T (L\bv_k)}{\bv_k^T \bv_k} = \frac{\bv_k^T L \bv_k}{\|\bv_k\|_2^2} \ge 0
\]
\item \textbf{Zero row sums ($L\11 = \b0$):} For each row $i = 1, \ldots, n$,
\[
(L\11)_i = \sum_{j=1}^n \ell_{ij} \cdot 1 = d_i - \sum_{j=1}^n a_{ij} = d_i - d_i = 0 \implies L\11 = \b0 = 0\cdot \11 \qquad \qed
\]
\end{enumerate}
\end{frame}


\begin{frame}{Beware of the eigenvalue ordering convention!}
In PCA (Unit 05) and Perron-Frobenius theory (Units 08--09), we ordered eigenvalues from \emph{largest to smallest} because we wanted maximum variance or the dominant long-term state
\vfill
For the graph Laplacian $L$, the universal convention is the \textbf{ascending order}
\[
0 = \lambda_1 \le \lambda_2 \le \cdots \le \lambda_n
\]
\vfill
Why? Because for any unit eigenvector $\bv_k$ ($\|\bv_k\|_2 = 1$),
\[
\lambda_k = \bv_k^T L \bv_k = \sum_{\{i, j\} \in E} a_{ij}(v_{ki} - v_{kj})^2
\]
\begin{itemize}
\item \textbf{Smallest eigenvalues ($\lambda_1 = 0, \lambda_2, \ldots$)}: eigenvectors vary \emph{least} across edges, revealing connected components and large-scale community structure
\item \textbf{Largest eigenvalues ($\ldots, \lambda_n$)}: eigenvectors oscillate rapidly across edges
\end{itemize}
\end{frame}


\begin{frame}{Multiplicity of $\lambda = 0$ and connected components}
We know $\lambda_1 = 0$ is always an eigenvalue of $L$. When is $\lambda_2 = 0$ as well?
\vfill
\begin{importanttheorem}[Multiplicity of $0$ equals the number of connected components]
\label{th:zero-eigenvalues-components}
Let $G = (V, E)$ be an undirected graph with nonnegative edge weights and Laplacian $L$. The multiplicity of $0$ as an eigenvalue of $L$ (that is, $\nullity(L) = \dim(\nullspace(L))$) equals the number $k$ of connected components $C_1, C_2, \ldots, C_k$ of $G$
\vskip0.2cm
Moreover, the component indicator vectors $\11_{C_1}, \11_{C_2}, \ldots, \11_{C_k} \in \{0, 1\}^n$, defined by
\[
(\11_{C_r})_i = \begin{cases}
1 & \text{if } i \in C_r \\
0 & \text{otherwise}
\end{cases}
\]
form an orthogonal basis of $\nullspace(L)$
\end{importanttheorem}
\end{frame}


\begin{frame}[red]{Proof of Theorem~\ref{th:zero-eigenvalues-components} -- Part 1: connected graph ($k = 1$)}
Suppose first that $G$ is connected ($k = 1$). We must show that $\nullspace(L) = \Span\{\11\}$, so $\dim(\nullspace(L)) = 1$
\vfill
\begin{enumerate}
\item We already know $L\11 = \b0$, so $\Span\{\11\} \subseteq \nullspace(L)$
\vfill
\item Conversely, let $\bx \in \nullspace(L)$, so $L\bx = \b0$. Taking the dot product with $\bx$ and applying Theorem~\ref{th:laplacian-quadratic-form} gives
\[
0 = \bx^T \b0 = \bx^T L \bx = \sum_{\{i, j\} \in E} a_{ij}(x_i - x_j)^2
\]
\item For every edge $\{i, j\} \in E$, we have $a_{ij} > 0$ and $(x_i - x_j)^2 \ge 0$. A sum of nonnegative terms equals $0$ if and only if every term is $0$:
\[
(x_i - x_j)^2 = 0 \implies x_i = x_j \quad \text{for every edge } \{i, j\} \in E
\]
\item Since $G$ is connected, any two vertices $u, v \in V$ are linked by a path $u = i_0, i_1, \ldots, i_m = v$ in $G$, so $x_u = x_{i_1} = \cdots = x_v$
\vfill
\item Hence all entries of $\bx$ are equal ($x_1 = \cdots = x_n = c$), i.e. $\bx = c\,\11 \in \Span\{\11\}$ \hfill $\qed$
\end{enumerate}
\end{frame}


\begin{frame}[red]{Proof of Theorem~\ref{th:zero-eigenvalues-components} -- Part 2: $k$ connected components ($k \ge 2$)}
Now suppose $G$ has $k$ connected components $C_1, \ldots, C_k$. Reorder the vertices so that those in $C_1$ come first, followed by $C_2$, \ldots, $C_k$
\vfill
\begin{enumerate}
\item Because there are no edges between different components, $A$, $D$ and $L$ are block-diagonal:
\[
L = \begin{pmatrix}
L_1 & \b0 & \cdots & \b0 \\
\b0 & L_2 & \cdots & \b0 \\
\vdots & \vdots & \ddots & \vdots \\
\b0 & \b0 & \cdots & L_k
\end{pmatrix}
\]
where each $L_r$ is the Laplacian of the connected subgraph on $C_r$
\item Writing $\bx = (\bx^{(1)}, \ldots, \bx^{(k)})^T$ according to the blocks, $L\bx = \b0 \iff L_r \bx^{(r)} = \b0$ for every $r = 1, \ldots, k$
\item By Part 1 applied to each connected component $C_r$, $L_r \bx^{(r)} = \b0 \iff \bx^{(r)} = c_r \11_{|C_r|}$, which in $\IR^n$ means
\[
\bx = c_1 \11_{C_1} + c_2 \11_{C_2} + \cdots + c_k \11_{C_k}
\]
\item Since $C_1, \ldots, C_k$ are disjoint, $\11_{C_r}^T \11_{C_s} = 0$ for $r \neq s$, so $\{\11_{C_1}, \ldots, \11_{C_k}\}$ is an orthogonal basis of $\nullspace(L)$ and $\dim(\nullspace(L)) = k$ \hfill $\qed$
\end{enumerate}
\end{frame}


\begin{frame}{Example: a disconnected graph on 4 vertices}
Consider $G = (V, E)$ on $n = 4$ vertices with only two edges $\{1, 2\}$ and $\{3, 4\}$, so $G$ has $k = 2$ connected components $C_1 = \{1, 2\}$ and $C_2 = \{3, 4\}$:
\vfill
\begin{columns}[c]
\begin{column}{0.32\textwidth}
\centering
\begin{tikzpicture}[scale=1.1, every node/.style={circle, draw=black, thick, fill=yellow!20, minimum size=7mm, font=\small\bfseries}]
\node[fill=blue!15] (1) at (0, 1.3) {1};
\node[fill=blue!15] (2) at (1.5, 1.3) {2};
\node[fill=red!15] (3) at (0, 0) {3};
\node[fill=red!15] (4) at (1.5, 0) {4};
\draw[very thick, blue!70!black] (1) -- (2);
\draw[very thick, red!70!black] (3) -- (4);
\end{tikzpicture}
\end{column}
\begin{column}{0.68\textwidth}
\[
L = \left(\begin{array}{cc|cc}
1 & -1 & 0 & 0 \\
-1 & 1 & 0 & 0 \\ \hline
0 & 0 & 1 & -1 \\
0 & 0 & -1 & 1
\end{array}\right) = \begin{pmatrix} L_1 & \b0 \\ \b0 & L_2 \end{pmatrix}
\]
\end{column}
\end{columns}
\vfill
\begin{itemize}
\item Each $2\times 2$ block $\begin{pmatrix} 1 & -1 \\ -1 & 1 \end{pmatrix}$ has trace $2$ and determinant $0$, so its eigenvalues are $0$ and $2$
\item The spectrum of $L$ is therefore $\lambda_1 = 0, \; \lambda_2 = 0, \; \lambda_3 = 2, \; \lambda_4 = 2$ (multiplicity of $0$ is $k = 2$)
\item The orthogonal indicator vectors for $\nullspace(L)$ are
\[
\11_{C_1} = (1, 1, 0, 0)^T \qquad \text{and} \qquad \11_{C_2} = (0, 0, 1, 1)^T
\]
\end{itemize}
\end{frame}


\begin{frame}[fragile]{Counting connected components with $\lambda = 0$ in R}
Let us build a graph on $n = 8$ vertices consisting of $k = 3$ disconnected components (a triangle $\{1, 2, 3\}$, a path $\{4, 5, 6\}$ and an edge $\{7, 8\}$) and check the eigenvalues of $L$:
\vfill
\begin{knitrout}
\definecolor{shadecolor}{rgb}{0.969, 0.969, 0.969}\color{fgcolor}\begin{kframe}
\begin{alltt}
\hldef{G_comp} \hlkwb{=} \hlkwd{make_graph}\hldef{(}\hlkwc{edges} \hldef{=} \hlkwd{c}\hldef{(}\hlnum{1}\hldef{,}\hlnum{2}\hldef{,} \hlnum{2}\hldef{,}\hlnum{3}\hldef{,} \hlnum{1}\hldef{,}\hlnum{3}\hldef{,}  \hlnum{4}\hldef{,}\hlnum{5}\hldef{,} \hlnum{5}\hldef{,}\hlnum{6}\hldef{,}  \hlnum{7}\hldef{,}\hlnum{8}\hldef{),}
                    \hlkwc{n} \hldef{=} \hlnum{8}\hldef{,} \hlkwc{directed} \hldef{=} \hlnum{FALSE}\hldef{)}
\hldef{L_comp} \hlkwb{=} \hlkwd{as.matrix}\hldef{(}\hlkwd{laplacian_matrix}\hldef{(G_comp))}
\hldef{ev_comp} \hlkwb{=} \hlkwd{eigen}\hldef{(L_comp,} \hlkwc{symmetric} \hldef{=} \hlnum{TRUE}\hldef{)}
\hlcom{# Eigenvalues of L in ascending order}
\hlkwd{round}\hldef{(}\hlkwd{rev}\hldef{(ev_comp}\hlopt{$}\hldef{values),} \hlnum{4}\hldef{)}
\end{alltt}
\begin{verbatim}
## [1] 0 0 0 1 2 3 3 3
\end{verbatim}
\begin{alltt}
\hlcom{# Number of zero eigenvalues (up to numerical precision)}
\hlkwd{sum}\hldef{(}\hlkwd{abs}\hldef{(ev_comp}\hlopt{$}\hldef{values)} \hlopt{<} \hlnum{1e-10}\hldef{)}
\end{alltt}
\begin{verbatim}
## [1] 3
\end{verbatim}
\end{kframe}
\end{knitrout}
\vfill
Exactly three eigenvalues of $L$ equal $0$, matching the $k = 3$ connected components of \code{G\_comp}
\end{frame}


\begin{frame}{Spectrum of extreme graphs: empty vs complete graph $K_n$}
Comparing two extreme graphs on $n$ vertices illustrates how Laplacian eigenvalues measure connectivity:
\vfill
\begin{columns}[t]
\begin{column}{0.46\textwidth}
\textbf{1. Empty graph $E_n$ ($|E| = 0$):}
\[
L = \b0_n
\]
\begin{itemize}
\item $n$ isolated vertices ($k = n$ components)
\item Spectrum of $L$:
\[
\lambda_1 = \lambda_2 = \cdots = \lambda_n = 0
\]
\end{itemize}
\end{column}
\begin{column}{0.52\textwidth}
\textbf{2. Complete graph $K_n$ (all pairs connected):}
\[
L = (n - 1)\II_n - (\IJ_n - \II_n) = n\II_n - \11\11^T
\]
\begin{itemize}
\item For $\bv_1 = \11$: $L\11 = n\11 - n\11 = \b0 \implies \lambda_1 = 0$
\item For any $\bx \perp \11$ ($\11^T \bx = 0$):
\[
L\bx = n\bx - \11(\11^T \bx) = n\bx
\]
\item Spectrum of $K_n$: $\lambda_1 = 0$ and $\lambda_2 = \cdots = \lambda_n = n$
\end{itemize}
\end{column}
\end{columns}
\vfill
For every simple graph $G$ on $n$ vertices, its eigenvalues lie between these two extremes: $0 = \lambda_1 \le \lambda_2 \le \cdots \le \lambda_n \le n$
\end{frame}


\begin{frame}{What if $G$ is ``almost'' disconnected?}
Suppose $G$ has two dense clusters $C_1$ and $C_2$ that are completely disconnected ($k = 2$):
\begin{itemize}
\item Then $\lambda_1 = 0$ and $\lambda_2 = 0$, and any vector in $\nullspace(L)$ is constant on $C_1$ and constant on $C_2$:
\[
\bv_1 = (1, \ldots, 1, \; 1, \ldots, 1)^T, \qquad \bv_2 = (1, \ldots, 1, \; -1, \ldots, -1)^T \quad (\text{if } |C_1| = |C_2|)
\]
\end{itemize}
\vfill
Now perturb $G$ slightly by adding a single weak bridge edge between $C_1$ and $C_2$:
\begin{itemize}
\item The graph is now connected ($k = 1$), so $\lambda_1 = 0$ and $\lambda_2 > 0$
\item Because eigenvalues and eigenvectors of symmetric matrices vary continuously under small perturbations, $\lambda_2$ will be \textbf{positive but very small} ($\lambda_2 \approx 0$)
\item The associated eigenvector $\bv_2 \perp \11$ will remain \textbf{nearly constant and positive on $C_1$} and \textbf{nearly constant and negative on $C_2$}!
\end{itemize}
\vfill
Next, we make this intuition completely rigorous via \textbf{graph cuts}
\end{frame}


%%%%%%%%%%%%%%%%%%%%%
%%%%%%%%%%%%%%%%%%%%%
\SSsection{Graph cuts, the Fiedler vector and spectral clustering}

\SSsubsection{Graph partitioning and cut size}

\begin{frame}{Partitions and the cut of a graph}
Let $G = (V, E)$ be an undirected graph on $n$ vertices with symmetric edge weights $a_{ij} \ge 0$
\vfill
A \defword{2-way partition} (or \defword{bipartition}) of $G$ splits $V$ into two disjoint non-empty subsets $S$ and $\bar{S} = V \setminus S$ ($S \cup \bar{S} = V$, $S \cap \bar{S} = \emptyset$)
\vfill
\begin{importantdefinition}[Cut of a partition]
The \defword{cut} between $S$ and $\bar{S}$ is the total weight of edges with one endpoint in $S$ and the other endpoint in $\bar{S}$:
\[
\mathsf{cut}(S, \bar{S}) = \sum_{i \in S, \; j \in \bar{S}} a_{ij}
\]
For an unweighted graph ($a_{ij} \in \{0, 1\}$), $\mathsf{cut}(S, \bar{S})$ is simply the number of edges crossing between $S$ and $\bar{S}$
\end{importantdefinition}
\end{frame}


\begin{frame}{Encoding a partition with a $\pm 1$ indicator vector}
To connect graph cuts to linear algebra, we represent any bipartition $(S, \bar{S})$ by a \defword{cut indicator vector} $\bx = (x_1, \ldots, x_n)^T \in \{-1, +1\}^n$:
\[
x_i = \begin{cases}
+1 & \text{if } i \in S \\
-1 & \text{if } i \in \bar{S}
\end{cases}
\]
\vfill
\begin{columns}[c]
\begin{column}{0.48\textwidth}
\centering
\begin{tikzpicture}[scale=1.05]
\node[circle, draw=blue!80!black, very thick, fill=blue!15, minimum size=7mm, font=\small\bfseries] (1) at (0, 1.3) {$+1$};
\node[circle, draw=blue!80!black, very thick, fill=blue!15, minimum size=7mm, font=\small\bfseries] (2) at (0, 0) {$+1$};
\node[circle, draw=blue!80!black, very thick, fill=blue!15, minimum size=7mm, font=\small\bfseries] (3) at (1.3, 0.65) {$+1$};
\node[circle, draw=red!80!black, very thick, fill=red!15, minimum size=7mm, font=\small\bfseries] (4) at (2.8, 0.65) {$-1$};
\node[circle, draw=red!80!black, very thick, fill=red!15, minimum size=7mm, font=\small\bfseries] (5) at (4.1, 1.3) {$-1$};
\node[circle, draw=red!80!black, very thick, fill=red!15, minimum size=7mm, font=\small\bfseries] (6) at (4.1, 0) {$-1$};
\draw[very thick] (1) -- (2) -- (3) -- (1);
\draw[very thick, red!80!black] (3) -- (4);
\draw[very thick] (4) -- (5) -- (6) -- (4);
\draw[dashed, very thick, red!80!black] (2.05, -0.35) -- (2.05, 1.65);
\node[below, font=\footnotesize, text=blue!80!black] at (0.5, -0.3) {$S = \{1, 2, 3\}$};
\node[below, font=\footnotesize, text=red!80!black] at (3.5, -0.3) {$\bar{S} = \{4, 5, 6\}$};
\end{tikzpicture}
\end{column}
\begin{column}{0.52\textwidth}
\begin{importantproposition}[Cut size as a quadratic form]
\label{prop:cut-quadratic-form}
For any partition $(S, \bar{S})$ of $V$ with indicator vector $\bx \in \{-1, +1\}^n$,
\[
\bx^T L \bx = 4\,\mathsf{cut}(S, \bar{S})
\]
or equivalently,
\[
\mathsf{cut}(S, \bar{S}) = \frac{1}{4}\bx^T L \bx
\]
\end{importantproposition}
\end{column}
\end{columns}
\end{frame}


\begin{frame}[red]{Proof of Proposition~\ref{prop:cut-quadratic-form} ($\bx^T L \bx = 4\,\mathsf{cut}(S, \bar{S})$)}
Let $\bx \in \{-1, +1\}^n$ be the cut indicator vector for $(S, \bar{S})$. For any edge $\{i, j\} \in E$, look at the squared difference $(x_i - x_j)^2$:
\vfill
\begin{itemize}
\item \textbf{Both endpoints in $S$} ($i \in S, j \in S$): $x_i = +1$ and $x_j = +1$, so
\[
(x_i - x_j)^2 = (1 - 1)^2 = 0
\]
\item \textbf{Both endpoints in $\bar{S}$} ($i \in \bar{S}, j \in \bar{S}$): $x_i = -1$ and $x_j = -1$, so
\[
(x_i - x_j)^2 = (-1 - (-1))^2 = 0
\]
\item \textbf{Endpoints in different sets} ($i \in S, j \in \bar{S}$ or $i \in \bar{S}, j \in S$): $x_i - x_j = \pm 2$, so
\[
(x_i - x_j)^2 = (\pm 2)^2 = 4
\]
\end{itemize}
\vfill
Substituting into Theorem~\ref{th:laplacian-quadratic-form} ($\bx^T L \bx = \sum_{\{i, j\} \in E} a_{ij}(x_i - x_j)^2$), only the edges crossing between $S$ and $\bar{S}$ survive, each multiplied by $4$:
\[
\bx^T L \bx = \sum_{i \in S, \; j \in \bar{S}} 4\,a_{ij} = 4 \sum_{i \in S, \; j \in \bar{S}} a_{ij} = 4\,\mathsf{cut}(S, \bar{S}) \qquad \qed
\]
\end{frame}


\begin{frame}{Example: comparing two cuts on the 4-vertex paw graph}
Recall the 4-vertex paw graph with Laplacian
\[
L = \begin{pmatrix}
2 & -1 & -1 & 0 \\
-1 & 2 & -1 & 0 \\
-1 & -1 & 3 & -1 \\
0 & 0 & -1 & 1
\end{pmatrix}
\]
Let us compare two candidate partitions using $\mathsf{cut}(S, \bar{S}) = \frac{1}{4}\bx^T L \bx$:
\vfill
\begin{columns}[t]
\begin{column}{0.48\textwidth}
\textbf{Partition A:} $S_1 = \{1, 2, 3\}$, $\bar{S}_1 = \{4\}$
\[
\bx^{(1)} = (1, 1, 1, -1)^T
\]
\[
L\bx^{(1)} = \begin{pmatrix} 0 \\ 0 \\ 2 \\ -2 \end{pmatrix} \implies (\bx^{(1)})^T L \bx^{(1)} = 4
\]
\[
\mathsf{cut}(S_1, \bar{S}_1) = \frac{4}{4} = 1 \quad (\text{only edge } \{3, 4\})
\]
\end{column}
\begin{column}{0.48\textwidth}
\textbf{Partition B:} $S_2 = \{1, 3\}$, $\bar{S}_2 = \{2, 4\}$
\[
\bx^{(2)} = (1, -1, 1, -1)^T
\]
\[
L\bx^{(2)} = \begin{pmatrix} 2 \\ -4 \\ 4 \\ -2 \end{pmatrix} \implies (\bx^{(2)})^T L \bx^{(2)} = 12
\]
\[
\mathsf{cut}(S_2, \bar{S}_2) = \frac{12}{4} = 3 \quad (\text{edges } \{1,2\}, \{2,3\}, \{3,4\})
\]
\end{column}
\end{columns}
\end{frame}


%%%%%%%%%%%%%%%%%%%%%
\SSsubsection{Spectral relaxation and the Fiedler vector}

\begin{frame}{Why unconstrained minimum cut is not enough}
Since $\mathsf{cut}(S, \bar{S}) = \frac{1}{4}\bx^T L \bx$, can we cluster a graph simply by minimising $\bx^T L \bx$ over $\bx \in \{-1, +1\}^n$?
\vfill
\begin{enumerate}
\item \textbf{Trivial all-in-one solution:} Without constraints, the minimum of $\bx^T L \bx$ over $\{-1, +1\}^n$ is $0$, achieved at $\bx = +\11$ ($S = V, \bar{S} = \emptyset$) and $\bx = -\11$ ($S = \emptyset, \bar{S} = V$)
\vfill
\item \textbf{Unbalanced ``whisker'' cuts:} Even if we forbid $\bx = \pm\11$, minimum cut often slices off a single degree-1 vertex (like vertex $4$ in the paw graph) rather than splitting the graph into two meaningful communities of comparable size
\end{enumerate}
\vfill
To find balanced communities, we must add a \defword{balance constraint} requiring $|S|$ and $|\bar{S}|$ to be equal (or comparable) in size
\end{frame}


\begin{frame}{Encoding balanced bisection algebraically}
Suppose $n$ is even and we want a \defword{balanced bisection} with $|S| = |\bar{S}| = n/2$
\vfill
For any cut indicator vector $\bx \in \{-1, +1\}^n$, notice two algebraic facts:
\vfill
\begin{enumerate}
\item \textbf{Balance equals orthogonality to $\11$:} Since $x_i = +1$ on $S$ and $x_i = -1$ on $\bar{S}$,
\[
\11^T \bx = \sum_{i=1}^n x_i = |S| - |\bar{S}|
\]
Therefore,
\[
|S| = |\bar{S}| \quad \Longleftrightarrow \quad \11^T \bx = 0 \quad \Longleftrightarrow \quad \bx \perp \11
\]
\item \textbf{Fixed Euclidean norm:} Every vector $\bx \in \{-1, +1\}^n$ has squared norm
\[
\|\bx\|_2^2 = \bx^T \bx = \sum_{i=1}^n x_i^2 = \sum_{i=1}^n (\pm 1)^2 = n
\]
\end{enumerate}
\end{frame}


\begin{frame}{General balanced cuts: RatioCut for unequal cluster sizes}
What if the two true clusters do not have the exact same size ($|S| \neq |\bar{S}|$)? We penalise small sets by minimising the \defword{RatioCut} (Hagen \& Kahng, 1992):
\[
\mathsf{RatioCut}(S, \bar{S}) = \mathsf{cut}(S, \bar{S})\left(\frac{1}{|S|} + \frac{1}{|\bar{S}|}\right) = \frac{n}{|S|\,|\bar{S}|}\,\mathsf{cut}(S, \bar{S})
\]
\vfill
Define the weighted indicator vector $\bx = (x_1, \ldots, x_n)^T$ by
\[
x_i = \begin{cases}
+\sqrt{|\bar{S}| / |S|} & \text{if } i \in S \\[2pt]
-\sqrt{|S| / |\bar{S}|} & \text{if } i \in \bar{S}
\end{cases}
\]
Then one checks directly that:
\begin{enumerate}
\item $\11^T \bx = |S|\sqrt{|\bar{S}|/|S|} - |\bar{S}|\sqrt{|S|/|\bar{S}|} = \sqrt{|S|\,|\bar{S}|} - \sqrt{|S|\,|\bar{S}|} = 0$ ($\bx \perp \11$)
\item $\|\bx\|_2^2 = |S|\frac{|\bar{S}|}{|S|} + |\bar{S}|\frac{|S|}{|\bar{S}|} = |\bar{S}| + |S| = n$, and $\bx^T L \bx = n\,\mathsf{RatioCut}(S, \bar{S})$
\end{enumerate}
\end{frame}


\begin{frame}{From NP-hard discrete cut to continuous spectral relaxation}
Both balanced minimum bisection and RatioCut lead to minimising $\bx^T L \bx$ subject to $\bx \perp \11$, $\|\bx\|_2^2 = n$ and a discrete two-value constraint on the entries $x_i$:
\[
\min_{\substack{\bx \in \{-1, +1\}^n \\ \bx^T \11 = 0}} \frac{1}{4}\bx^T L \bx
\]
\vfill
\begin{itemize}
\item \textbf{Computational barrier:} Checking all balanced partitions $\binom{n}{n/2}$ is \textbf{NP-hard} (for $n = 100$ vertices, $\binom{100}{50} \approx 1.0 \times 10^{29}$)
\vfill
\item \textbf{Spectral relaxation:} Relax the discrete constraint on $x_i$ to allow \emph{arbitrary real values} $\bx \in \IR^n$ that still satisfy the norm constraint $\bx^T \bx = n$ and the balance constraint $\bx^T \11 = 0$:
\[
\min_{\substack{\bx \in \IR^n \\ \bx^T \bx = n, \;\; \bx^T \11 = 0}} \bx^T L \bx \quad \Longleftrightarrow \quad \min_{\substack{\bx \in \IR^n \setminus \{\b0\} \\ \bx \perp \11}} \frac{\bx^T L \bx}{\bx^T \bx}
\]
\end{itemize}
\end{frame}


\begin{frame}{Solving the relaxed problem via the Rayleigh quotient}
\begin{importanttheorem}[Rayleigh quotient theorem for $\lambda_2$]
\label{th:rayleigh-fiedler}
Let $L$ be the Laplacian of an undirected graph $G$ on $n$ vertices with eigenvalues $0 = \lambda_1 \le \lambda_2 \le \cdots \le \lambda_n$ and orthonormal eigenvectors $\bv_1 = \frac{1}{\sqrt{n}}\11, \bv_2, \ldots, \bv_n$. Then
\[
\min_{\substack{\bx \in \IR^n \setminus \{\b0\} \\ \bx \perp \11}} \frac{\bx^T L \bx}{\bx^T \bx} = \lambda_2
\]
and the minimum is attained when $\bx$ is any non-zero scalar multiple of $\bv_2$
\end{importanttheorem}
\vfill
In particular, if we scale $\bv_2$ so that $\|\bx\|_2^2 = n$, the optimal continuous vector is $\bx = \sqrt{n}\,\bv_2$, with objective value $\bx^T L \bx = n\lambda_2$ (giving the lower bound $\mathsf{cut}(S, \bar{S}) \ge \frac{n}{4}\lambda_2$ for any balanced bisection!)
\end{frame}


\begin{frame}[red]{Proof of Theorem~\ref{th:rayleigh-fiedler}}
Let $\{\bv_1, \bv_2, \ldots, \bv_n\}$ be an orthonormal basis of eigenvectors of $L$ with $L\bv_k = \lambda_k \bv_k$, where $0 = \lambda_1 \le \lambda_2 \le \cdots \le \lambda_n$ and $\bv_1 = \frac{1}{\sqrt{n}}\11$
\vfill
\begin{enumerate}
\item Expand any vector $\bx \in \IR^n \setminus \{\b0\}$ in the orthonormal basis $\{\bv_1, \ldots, \bv_n\}$:
\[
\bx = c_1 \bv_1 + c_2 \bv_2 + \cdots + c_n \bv_n, \qquad c_k = \bv_k^T \bx
\]
\item The orthogonality constraint $\bx \perp \11$ means $c_1 = \bv_1^T \bx = \frac{1}{\sqrt{n}}\11^T \bx = 0$, so
\[
\bx = \sum_{k=2}^n c_k \bv_k \qquad \text{and} \qquad \bx^T \bx = \sum_{k=2}^n c_k^2 > 0
\]
\item Since $L\bx = \sum_{k=2}^n c_k \lambda_k \bv_k$ and the $\bv_k$ are orthonormal,
\[
\frac{\bx^T L \bx}{\bx^T \bx} = \frac{\sum_{k=2}^n \lambda_k c_k^2}{\sum_{k=2}^n c_k^2} \ge \frac{\sum_{k=2}^n \lambda_2 c_k^2}{\sum_{k=2}^n c_k^2} = \lambda_2
\]
with equality when $c_2 \neq 0$ and $c_3 = \cdots = c_n = 0$, i.e. $\bx = c_2 \bv_2$ \hfill $\qed$
\end{enumerate}
\end{frame}


\begin{frame}{Algebraic connectivity and the Fiedler vector}
Miroslav Fiedler (1973) discovered the fundamental role of $\lambda_2$ and $\bv_2$ in measuring how well-connected a graph is
\vfill
\begin{importantdefinition}[Algebraic connectivity and Fiedler vector]
Let $G = (V, E)$ be an undirected graph on $n$ vertices with Laplacian eigenvalues $0 = \lambda_1 \le \lambda_2 \le \cdots \le \lambda_n$
\begin{itemize}
\item The second smallest eigenvalue $\lambda_2$ is called the \defword{algebraic connectivity} (or \defword{Fiedler value}) of $G$
\item An eigenvector $\bv_2$ associated with $\lambda_2$ (orthogonal to $\11$) is called a \defword{Fiedler vector} of $G$
\end{itemize}
\end{importantdefinition}
\vfill
\begin{itemize}
\item By Theorem~\ref{th:zero-eigenvalues-components}, $\lambda_2 > 0$ \textbf{if and only if $G$ is connected}
\item Because $\bv_2 \perp \11$, the entries of $\bv_2$ sum to $0$ ($\sum_{i=1}^n v_{2i} = 0$), so $\bv_2$ always has \textbf{both positive and negative entries}!
\end{itemize}
\end{frame}


\begin{frame}{The 2-way spectral partitioning algorithm}
How do we turn the continuous Fiedler vector $\bv_2 = (v_{21}, \ldots, v_{2n})^T \in \IR^n$ back into a discrete partition $(S, \bar{S})$? We round its entries by their signs!
\vfill
\begin{block}{Algorithm: 2-way unnormalised spectral clustering (Fiedler bisection)}
\textbf{Input:} Undirected graph $G = (V, E)$ with symmetric adjacency/similarity matrix $A \in \M_n(\IRplus)$
\begin{enumerate}
\item Compute the degree matrix $D = \diag(A\11)$ and the Laplacian $L = D - A$
\item Compute the second smallest eigenvalue $\lambda_2$ and associated Fiedler vector $\bv_2 = (v_{21}, \ldots, v_{2n})^T$ of $L$
\item Partition the vertices according to the sign of $v_{2i}$:
\[
S = \{i \in V : v_{2i} \ge 0\}, \qquad \bar{S} = \{i \in V : v_{2i} < 0\}
\]
(Alternatively, split at the median of $\bv_2$ if an exact $n/2 - n/2$ split is required)
\end{enumerate}
\end{block}
\vfill
Note: Because $-\bv_2$ is also a Fiedler vector, flipping the sign of $\bv_2$ simply swaps the labels $S \leftrightarrow \bar{S}$, which defines the \emph{exact same} partition of $V$
\end{frame}


\begin{frame}{Extension to $k$ clusters and normalised Laplacians}
What if a graph (or dataset) has $k > 2$ clusters, or vertices with wildly different degrees?
\vfill
\begin{itemize}
\item \textbf{$k$-way spectral clustering:}
\begin{enumerate}
\item Compute the first $k$ eigenvectors $\bv_1, \bv_2, \ldots, \bv_k$ of $L$ associated with the $k$ smallest eigenvalues $0 = \lambda_1 \le \lambda_2 \le \cdots \le \lambda_k$
\item Stack them as columns of the matrix $U = [\bv_1 \;\; \bv_2 \;\; \cdots \;\; \bv_k] \in \M_{n, k}(\IR)$
\item Treat row $i$ of $U$ as the new embedding $\by_i \in \IR^k$ of vertex $i$, and run $k$-means on $\{\by_1, \ldots, \by_n\} \subset \IR^k$
\end{enumerate}
\vfill
\item \textbf{Normalised graph Laplacians} (used when degrees $d_i$ vary over orders of magnitude):
\[
L_{\text{sym}} = D^{-1/2} L D^{-1/2} = \II - D^{-1/2} A D^{-1/2}, \qquad
L_{\text{rw}} = D^{-1} L = \II - D^{-1} A
\]
Notice that $D^{-1} A$ is the random-walk transition matrix on $G$! Spectral clustering with $L_{\text{sym}}$ or $L_{\text{rw}}$ relaxes the degree-normalised cut (Shi \& Malik 2000; Ng, Jordan \& Weiss 2001)
\end{itemize}
\end{frame}


\begin{frame}{Choosing the number of clusters $k$: the eigengap heuristic}
How do we choose the number of clusters $k$ from the spectrum $0 = \lambda_1 \le \lambda_2 \le \cdots \le \lambda_n$?
\vfill
\begin{itemize}
\item By Theorem~\ref{th:zero-eigenvalues-components}, if $G$ consists of $k$ completely disconnected clusters, then
\[
\lambda_1 = \lambda_2 = \cdots = \lambda_k = 0 \qquad \text{and} \qquad \lambda_{k+1} > 0
\]
\item When those $k$ dense clusters are joined by only a few weak between-cluster edges, the first $k$ eigenvalues remain \textbf{very close to $0$}, while $\lambda_{k+1}$ jumps sharply upward
\vfill
\item \defword{Eigengap heuristic}: plot $\lambda_1, \lambda_2, \ldots, \lambda_m$ in ascending order and look for the index $k$ with the largest \textbf{eigengap}
\[
\Delta_k = \lambda_{k+1} - \lambda_k
\]
\end{itemize}
\vfill
For example, if $\lambda_1 = 0, \; \lambda_2 = 0.04, \; \lambda_3 = 0.07, \; \lambda_4 = 2.85, \; \lambda_5 = 3.10$, the large gap $\Delta_3 = \lambda_4 - \lambda_3 = 2.78$ strongly indicates $k = 3$ clusters
\end{frame}


%%%%%%%%%%%%%%%%%%%%%
\SSsubsection{Worked exam-style examples and R implementation}

\begin{frame}{Example 1: spectral bisection of the 4-vertex path $P_4$}
Consider the 4-vertex path graph $P_4$:
\begin{center}
\begin{tikzpicture}[scale=1.15, every node/.style={circle, draw=black, thick, fill=yellow!20, minimum size=7mm, font=\small\bfseries}]
\node (1) at (0, 0) {1};
\node (2) at (1.8, 0) {2};
\node (3) at (3.6, 0) {3};
\node (4) at (5.4, 0) {4};
\draw[very thick] (1) -- (2) -- (3) -- (4);
\end{tikzpicture}
\end{center}
\vfill
\textbf{Step 1 -- Degree, adjacency and Laplacian matrices:}
\[
D = \begin{pmatrix}
1 & 0 & 0 & 0 \\
0 & 2 & 0 & 0 \\
0 & 0 & 2 & 0 \\
0 & 0 & 0 & 1
\end{pmatrix}, \quad
A = \begin{pmatrix}
0 & 1 & 0 & 0 \\
1 & 0 & 1 & 0 \\
0 & 1 & 0 & 1 \\
0 & 0 & 1 & 0
\end{pmatrix} \implies
L = D - A = \begin{pmatrix}
1 & -1 & 0 & 0 \\
-1 & 2 & -1 & 0 \\
0 & -1 & 2 & -1 \\
0 & 0 & -1 & 1
\end{pmatrix}
\]
\end{frame}


\begin{frame}{Example 1 (continued): Fiedler vector and cut of $P_4$}
\textbf{Step 2 -- Verify the Fiedler eigenpair $(\lambda_2, \bv_2)$:}
Claim: $\lambda_2 = 2 - \sqrt{2} \approx 0.586$ and $\bv_2 = \begin{pmatrix} 1, & \sqrt{2}-1, & 1-\sqrt{2}, & -1 \end{pmatrix}^T \approx \begin{pmatrix} 1, & 0.414, & -0.414, & -1 \end{pmatrix}^T$
\vfill
Let us verify $L\bv_2 = \lambda_2 \bv_2$ directly:
\[
L\bv_2 = \begin{pmatrix}
1 & -1 & 0 & 0 \\
-1 & 2 & -1 & 0 \\
0 & -1 & 2 & -1 \\
0 & 0 & -1 & 1
\end{pmatrix}
\begin{pmatrix}
1 \\ \sqrt{2}-1 \\ 1-\sqrt{2} \\ -1
\end{pmatrix}
= \begin{pmatrix}
2 - \sqrt{2} \\
3\sqrt{2} - 4 \\
4 - 3\sqrt{2} \\
\sqrt{2} - 2
\end{pmatrix}
= (2 - \sqrt{2})
\begin{pmatrix}
1 \\ \sqrt{2}-1 \\ 1-\sqrt{2} \\ -1
\end{pmatrix} = \lambda_2 \bv_2
\]
(since $(2-\sqrt{2})(\sqrt{2}-1) = 3\sqrt{2}-4$). Notice also that $\11^T \bv_2 = 1 + (\sqrt{2}-1) + (1-\sqrt{2}) - 1 = 0$
\vfill
\textbf{Step 3 -- Partition by the signs of $\bv_2$:}
\begin{itemize}
\item $v_{21} = 1 > 0$ and $v_{22} = \sqrt{2}-1 \approx 0.414 > 0 \implies S = \{1, 2\}$
\item $v_{23} = 1-\sqrt{2} \approx -0.414 < 0$ and $v_{24} = -1 < 0 \implies \bar{S} = \{3, 4\}$
\item The spectral cut severing the middle edge $\{2, 3\}$ has $\mathsf{cut}(S, \bar{S}) = 1$
\end{itemize}
\end{frame}


\begin{frame}{Example 2: the 6-vertex barbell graph}
Consider two triangles $C_1 = \{1, 2, 3\}$ and $C_2 = \{4, 5, 6\}$ joined by a single bridge edge $\{3, 4\}$:
\vfill
\begin{columns}[c]
\begin{column}{0.44\textwidth}
\centering
\begin{tikzpicture}[scale=1.05]
\node[circle, draw=black, thick, fill=blue!15, minimum size=7mm, font=\small\bfseries] (1) at (0, 1.3) {1};
\node[circle, draw=black, thick, fill=blue!15, minimum size=7mm, font=\small\bfseries] (2) at (0, 0) {2};
\node[circle, draw=black, thick, fill=blue!15, minimum size=7mm, font=\small\bfseries] (3) at (1.25, 0.65) {3};
\node[circle, draw=black, thick, fill=red!15, minimum size=7mm, font=\small\bfseries] (4) at (2.65, 0.65) {4};
\node[circle, draw=black, thick, fill=red!15, minimum size=7mm, font=\small\bfseries] (5) at (3.9, 1.3) {5};
\node[circle, draw=black, thick, fill=red!15, minimum size=7mm, font=\small\bfseries] (6) at (3.9, 0) {6};
\draw[very thick] (1) -- (2) -- (3) -- (1);
\draw[very thick, red!80!black] (3) -- (4);
\draw[very thick] (4) -- (5) -- (6) -- (4);
\end{tikzpicture}
\vskip0.2cm
\footnotesize
Degrees: $d_1 = d_2 = d_5 = d_6 = 2$, $d_3 = d_4 = 3$
\end{column}
\begin{column}{0.56\textwidth}
\[
L = \left(\begin{array}{ccc|ccc}
2 & -1 & -1 & 0 & 0 & 0 \\
-1 & 2 & -1 & 0 & 0 & 0 \\
-1 & -1 & 3 & -1 & 0 & 0 \\ \hline
0 & 0 & -1 & 3 & -1 & -1 \\
0 & 0 & 0 & -1 & 2 & -1 \\
0 & 0 & 0 & -1 & -1 & 2
\end{array}\right)
\]
\end{column}
\end{columns}
\vfill
Notice that $L$ is nearly block-diagonal: the two $3\times 3$ blocks are coupled only through $\ell_{34} = \ell_{43} = -1$ and the extra $+1$ on $\ell_{33}$ and $\ell_{44}$
\end{frame}


\begin{frame}{Example 2 (continued): exact Fiedler vector of the barbell}
By symmetry between $\{1, 2\}$ and $\{5, 6\}$ and antisymmetry across the bridge $\{3, 4\}$ (so that $\11^T \bv_2 = 0$), we look for a Fiedler vector of the form
\[
\bv_2 = (1, \; 1, \; c, \; -c, \; -1, \; -1)^T
\]
\vfill
Substitute into $L\bv_2 = \lambda \bv_2$:
\begin{itemize}
\item \textbf{Row 1 (and Row 2):} $2(1) - 1(1) - 1(c) = 1 - c = \lambda \implies c = 1 - \lambda$
\item \textbf{Row 3:} $-1(1) - 1(1) + 3c - 1(-c) = 4c - 2 = \lambda c$
\end{itemize}
\vfill
Substitute $c = 1 - \lambda$ into $4c - 2 = \lambda c$:
\[
4(1 - \lambda) - 2 = \lambda(1 - \lambda) \iff \lambda^2 - 5\lambda + 2 = 0 \implies \lambda_2 = \frac{5 - \sqrt{17}}{2} \approx 0.4384
\]
Then $c = 1 - \lambda_2 = \frac{\sqrt{17} - 3}{2} \approx 0.5616 > 0$, giving the Fiedler vector
\[
\bv_2 \approx (1, \; 1, \; 0.562, \; -0.562, \; -1, \; -1)^T
\]
Splitting by sign ($\ge 0$ vs $< 0$) gives $S = \{1, 2, 3\}$ and $\bar{S} = \{4, 5, 6\}$, cutting only the bridge $\{3, 4\}$!
\end{frame}


\begin{frame}[fragile]{Computing the barbell Fiedler vector in R}
We can build the barbell Laplacian in R (either from $D - A$ or with \code{igraph::laplacian\_matrix}) and inspect its eigenvalues and Fiedler vector:
\vfill
\begin{knitrout}
\definecolor{shadecolor}{rgb}{0.969, 0.969, 0.969}\color{fgcolor}\begin{kframe}
\begin{alltt}
\hldef{A_barbell} \hlkwb{=} \hlkwd{matrix}\hldef{(}\hlkwd{c}\hldef{(}
  \hlnum{0}\hldef{,} \hlnum{1}\hldef{,} \hlnum{1}\hldef{,} \hlnum{0}\hldef{,} \hlnum{0}\hldef{,} \hlnum{0}\hldef{,}
  \hlnum{1}\hldef{,} \hlnum{0}\hldef{,} \hlnum{1}\hldef{,} \hlnum{0}\hldef{,} \hlnum{0}\hldef{,} \hlnum{0}\hldef{,}
  \hlnum{1}\hldef{,} \hlnum{1}\hldef{,} \hlnum{0}\hldef{,} \hlnum{1}\hldef{,} \hlnum{0}\hldef{,} \hlnum{0}\hldef{,}
  \hlnum{0}\hldef{,} \hlnum{0}\hldef{,} \hlnum{1}\hldef{,} \hlnum{0}\hldef{,} \hlnum{1}\hldef{,} \hlnum{1}\hldef{,}
  \hlnum{0}\hldef{,} \hlnum{0}\hldef{,} \hlnum{0}\hldef{,} \hlnum{1}\hldef{,} \hlnum{0}\hldef{,} \hlnum{1}\hldef{,}
  \hlnum{0}\hldef{,} \hlnum{0}\hldef{,} \hlnum{0}\hldef{,} \hlnum{1}\hldef{,} \hlnum{1}\hldef{,} \hlnum{0}
\hldef{),} \hlkwc{nrow} \hldef{=} \hlnum{6}\hldef{,} \hlkwc{byrow} \hldef{=} \hlnum{TRUE}\hldef{)}
\hldef{L_barbell} \hlkwb{=} \hlkwd{diag}\hldef{(}\hlkwd{rowSums}\hldef{(A_barbell))} \hlopt{-} \hldef{A_barbell}
\hldef{ev_b} \hlkwb{=} \hlkwd{eigen}\hldef{(L_barbell,} \hlkwc{symmetric} \hldef{=} \hlnum{TRUE}\hldef{)}
\hlkwd{round}\hldef{(}\hlkwd{rev}\hldef{(ev_b}\hlopt{$}\hldef{values),} \hlnum{4}\hldef{)}          \hlcom{# eigenvalues in ascending order}
\end{alltt}
\begin{verbatim}
## [1] 0.0000 0.4384 3.0000 3.0000 3.0000 4.5616
\end{verbatim}
\begin{alltt}
\hlkwd{round}\hldef{(ev_b}\hlopt{$}\hldef{vectors[,} \hlnum{5}\hldef{],} \hlnum{4}\hldef{)}         \hlcom{# unit Fiedler vector v_2}
\end{alltt}
\begin{verbatim}
## [1]  0.4647  0.4647  0.2610 -0.2610 -0.4647 -0.4647
\end{verbatim}
\end{kframe}
\end{knitrout}
\end{frame}


\begin{frame}[fragile]{Real-world benchmark: Zachary's Karate club network}
Zachary (1977) recorded social interactions among $n = 34$ members of a university karate club just before a dispute between the instructor (Mr Hi, vertex $1$) and the club president (John A, vertex $34$) split the club into two factions
\vfill
\begin{knitrout}
\definecolor{shadecolor}{rgb}{0.969, 0.969, 0.969}\color{fgcolor}\begin{kframe}
\begin{alltt}
\hldef{G_karate} \hlkwb{=} \hlkwd{make_graph}\hldef{(}\hlsng{"Zachary"}\hldef{)}
\hldef{L_karate} \hlkwb{=} \hlkwd{as.matrix}\hldef{(}\hlkwd{laplacian_matrix}\hldef{(G_karate))}
\hldef{ev_k} \hlkwb{=} \hlkwd{eigen}\hldef{(L_karate,} \hlkwc{symmetric} \hldef{=} \hlnum{TRUE}\hldef{)}
\hldef{n_k} \hlkwb{=} \hlkwd{vcount}\hldef{(G_karate)}
\hlcom{# Smallest three eigenvalues of L}
\hlkwd{round}\hldef{(ev_k}\hlopt{$}\hldef{values[n_k}\hlopt{:}\hldef{(n_k} \hlopt{-} \hlnum{2}\hldef{)],} \hlnum{4}\hldef{)}
\end{alltt}
\begin{verbatim}
## [1] 0.0000 0.4685 0.9092
\end{verbatim}
\begin{alltt}
\hlcom{# Fiedler vector v_2 (oriented so vertex 1 is positive)}
\hldef{v2_k} \hlkwb{=} \hldef{ev_k}\hlopt{$}\hldef{vectors[, n_k} \hlopt{-} \hlnum{1}\hldef{]}
\hlkwa{if} \hldef{(v2_k[}\hlnum{1}\hldef{]} \hlopt{<} \hlnum{0}\hldef{) v2_k} \hlkwb{=} \hlopt{-}\hldef{v2_k}
\hlkwd{table}\hldef{(}\hlkwc{Sign_of_v2} \hldef{=} \hlkwd{ifelse}\hldef{(v2_k} \hlopt{>=} \hlnum{0}\hldef{,} \hlsng{"Group 1 (+)"}\hldef{,} \hlsng{"Group 2 (-)"}\hldef{))}
\end{alltt}
\begin{verbatim}
## Sign_of_v2
## Group 1 (+) Group 2 (-) 
##          15          19
\end{verbatim}
\end{kframe}
\end{knitrout}
\end{frame}




\maxFrameImage{FIGS/U10-plot-karate-fiedler-1.pdf}


\begin{frame}[fragile]{$k$-way spectral clustering in R ($k = 3$ communities)}
Suppose a network has $k = 3$ dense communities of $15$ vertices each ($n = 45$), with high within-community edge probability ($p_{\text{in}} = 0.7$) and low between-community probability ($p_{\text{out}} = 0.03$):
\vfill
\begin{knitrout}
\definecolor{shadecolor}{rgb}{0.969, 0.969, 0.969}\color{fgcolor}\begin{kframe}
\begin{alltt}
\hlkwd{set.seed}\hldef{(}\hlnum{2740}\hldef{)}
\hldef{pm} \hlkwb{=} \hlkwd{matrix}\hldef{(}\hlnum{0.03}\hldef{,} \hlkwc{nrow} \hldef{=} \hlnum{3}\hldef{,} \hlkwc{ncol} \hldef{=} \hlnum{3}\hldef{);} \hlkwd{diag}\hldef{(pm)} \hlkwb{=} \hlnum{0.70}
\hldef{G_sbm} \hlkwb{=} \hlkwd{sample_sbm}\hldef{(}\hlnum{45}\hldef{,} \hlkwc{pref.matrix} \hldef{= pm,} \hlkwc{block.sizes} \hldef{=} \hlkwd{c}\hldef{(}\hlnum{15}\hldef{,} \hlnum{15}\hldef{,} \hlnum{15}\hldef{))}
\hldef{L_sbm} \hlkwb{=} \hlkwd{as.matrix}\hldef{(}\hlkwd{laplacian_matrix}\hldef{(G_sbm))}
\hldef{ev_sbm} \hlkwb{=} \hlkwd{eigen}\hldef{(L_sbm,} \hlkwc{symmetric} \hldef{=} \hlnum{TRUE}\hldef{)}
\hlcom{# Smallest 6 eigenvalues of L (note the jump after lambda_3!)}
\hlkwd{round}\hldef{(}\hlkwd{rev}\hldef{(ev_sbm}\hlopt{$}\hldef{values)[}\hlnum{1}\hlopt{:}\hlnum{6}\hldef{],} \hlnum{3}\hldef{)}
\end{alltt}
\begin{verbatim}
## [1] 0.000 1.203 1.500 5.418 6.865 7.517
\end{verbatim}
\begin{alltt}
\hlcom{# Spectral embedding using v_2 and v_3, followed by k-means (k = 3)}
\hldef{U_sbm} \hlkwb{=} \hldef{ev_sbm}\hlopt{$}\hldef{vectors[,} \hlkwd{c}\hldef{(}\hlnum{44}\hldef{,} \hlnum{43}\hldef{)]}
\hldef{km_sbm} \hlkwb{=} \hlkwd{kmeans}\hldef{(U_sbm,} \hlkwc{centers} \hldef{=} \hlnum{3}\hldef{,} \hlkwc{nstart} \hldef{=} \hlnum{20}\hldef{)}
\hlkwd{table}\hldef{(}\hlkwc{True_block} \hldef{=} \hlkwd{rep}\hldef{(}\hlnum{1}\hlopt{:}\hlnum{3}\hldef{,} \hlkwc{each} \hldef{=} \hlnum{15}\hldef{),} \hlkwc{Spectral_cluster} \hldef{= km_sbm}\hlopt{$}\hldef{cluster)}
\end{alltt}
\begin{verbatim}
##           Spectral_cluster
## True_block  1  2  3
##          1  0  0 15
##          2 15  0  0
##          3  0 15  0
\end{verbatim}
\end{kframe}
\end{knitrout}
\end{frame}




\maxFrameImage{FIGS/U10-plot-three-way-spectral-1.pdf}


\begin{frame}{Summary of Unit 10}
\begin{itemize}
\item For an undirected graph $G = (V, E)$ with symmetric adjacency/weight matrix $A$ and degree matrix $D = \diag(A\11)$, the \textbf{combinatorial Laplacian} is $L = D - A = B^T B$
\vfill
\item The \textbf{Laplacian quadratic form} measures smoothness across edges:
\[
\bx^T L \bx = \frac{1}{2}\sum_{i=1}^n \sum_{j=1}^n a_{ij}(x_i - x_j)^2 = \sum_{\{i, j\} \in E} a_{ij}(x_i - x_j)^2 \ge 0
\]
\vfill
\item $L$ is real, symmetric and positive semi-definite with eigenvalues $0 = \lambda_1 \le \lambda_2 \le \cdots \le \lambda_n$ and $L\11 = \b0$; the multiplicity of $\lambda = 0$ equals the number of connected components of $G$
\vfill
\item For any $\pm 1$ cut indicator vector $\bx \in \{-1, +1\}^n$, $\mathsf{cut}(S, \bar{S}) = \frac{1}{4}\bx^T L \bx$
\vfill
\item Relaxing balanced minimum cut ($\bx \perp \11$, $\|\bx\|_2^2 = n$) to $\IR^n$ yields the \textbf{Fiedler value} $\lambda_2$ and \textbf{Fiedler vector} $\bv_2$, whose signs partition $G$ into two communities
\end{itemize}
\end{frame}




% Save counters for next file

% Save theorem count for next file
\end{document}
